%% Follow comments to support use.

%%%%%%%%%%%%%%%%%%%%%%%%%%%%%%%%%%%%%%%%%%%%%%%%%%%%%%%%%
%% STEP 1: Choose options for MSc / BSc / seminar layout and your bibliographic style
%%%%%%%%%%%%%%%%%%%%%%%%%%%%%%%%%%%%%%%%%%%%%%%%%%%%%%%%%

%%  Language: 
%%      finnish, swedish, or english
%%  Pagination (use twoside by default)  
%%      oneside or twoside,
%%  Study programme / kind of report
%%      csm  = Master's thesis in Computer Science Master's Programme;
%%      tkt = Bachelor's thesis in Computer Science Bachelor's Programme;
%%      seminar = seminar report
%%  For Master's thesis choose your line or track:
%%      (30 cr thesis, 2020 onwards, Master's Programme in Computer Science = csm)
%%      software-track-2020 = Software study track
%%      algorithms-track-2020 = Algorithms study track
%%      networking-track-2020 = Networking study track

\documentclass[english,twoside,censored,csm,software-track-2020]{HYthesisML} 


% If wanted, open new chapters only at right page.
% By default, "openany".
%\PassOptionsToClass{openright,twoside,a4paper}{report}
\PassOptionsToClass{openany,twoside,a4paper}{report}

\usepackage{csquotes}
%%%%%%%%%%%%%%%%%%%%%%%%%%%%%%%%%%%%%%%%%%%%%%%%%%%%%%%%%
%% REFERENCES
%% Some notes on bibliography usage and options:
%% natbib -> you can use, e.g., \citep{} or \parencite{} for (Einstein, 1905); with APA \cite -> Einstein, 1905 without ()
%% maxcitenames=2 -> only 2 author names in text citations, if more -> et al. is used
%% maxbibnames=99 as no great need to suppress the biliography list in a thesis
%% for more information see biblatex package documentation, e.g., from https://ctan.org/pkg/biblatex 

%% Reference style: select one 
%% for APA = Harvard style = authoryear -> (Einstein, 1905) use:
\usepackage[style=authoryear,bibstyle=authoryear,backend=biber,natbib=true,maxnames=99,maxcitenames=2,uniquelist=minyear,giveninits=true,uniquename=mininit,urldate=long]{biblatex}
%% for numeric = Vancouver style -> [1] use:
%\usepackage[style=numeric,bibstyle=numeric,backend=biber,natbib=true,maxbibnames=99,giveninits=true,uniquename=init]{biblatex}
%% for alpahbetic -> [Ein05] use:
%\usepackage[style=alphabetic,bibstyle=alphabetic,backend=biber,natbib=true,maxbibnames=99,giveninits=true,uniquename=init]{biblatex}
%

\addbibresource{bibliography.bib}
% in case you want the final delimiter between authors & -> (Einstein & Zweistein, 1905) 
% \renewcommand{\finalnamedelim}{ \& }
% List the authors in the Bibilipgraphy as Lastname F, Familyname G,
\DeclareNameAlias{sortname}{family-given}
% remove the punctuation between author names in Bibliography 
%\renewcommand{\revsdnamepunct}{ }


%% Block of definitions for fonts and packages for picture management.
%% In some systems, the figure packages may not be happy together.
%% Choose the ones you need.

%\usepackage[utf8]{inputenc} % For UTF8 support, in some systems. Use UTF8 when saving your file.

\usepackage{lmodern}         % Font package, again in some systems.
\usepackage{textcomp}        % Package for special symbols
\usepackage[pdftex]{color, graphicx} % For pdf output and jpg/png graphics
\usepackage{epsfig}
\usepackage{subfigure}
\usepackage[pdftex, plainpages=false]{hyperref} % For hyperlinks and pdf metadata
\usepackage{fancyhdr}        % For nicer page headers
\usepackage{tikz}            % For making vector graphics (hard to learn but powerful)
%\usepackage{wrapfig}        % For nice text-wrapping figures (use at own discretion)
\usepackage{amsmath, amssymb} % For better math

\usepackage[normalem]{ulem} % For table
\useunder{\uline}{\ul}{} % For table
\usepackage{multirow} % For table
\usepackage{adjustbox} % For table

\singlespacing               %line spacing options; normally use single

\fussy
%\sloppy                      % sloppy and fussy commands can be used to avoid overlong text lines
% if you want to see which lines are too long or have too little stuff, comment out the following lines
% \overfullrule=1mm
% to see more info in the detailed log about under/overfull boxes...
% \showboxbreadth=50 
% \showboxdepth=50



%%%%%%%%%%%%%%%%%%%%%%%%%%%%%%%%%%%%%%%%%%%%%%%%%%%%%%%%%
%% STEP 2:
%%%%%%%%%%%%%%%%%%%%%%%%%%%%%%%%%%%%%%%%%%%%%%%%%%%%%%%%%
%% Set up personal information for the title page and the abstract form.
%% Replace parameters with your information.
\title{Capabilities of Generative AI in Web Accessibility Evaluation: A Design Science Approach}

\author{Sebastian Sergelius}
\date{\today}

% Set supervisors, use the titles according to the thesis language
% in English Prof. or Dr., or in Finnish toht. or tri or FT, TkT, Ph.D. or in Swedish... 
\supervisors{Prof. Tomi Männistö}

\keywords{algorithms, data structures}
\additionalinformation{\translate{\track}}

%% For seminar reports:
%%\additionalinformation{Name of the seminar}

%% Provide classification terms, to appear on the abstract page.
%% Replace the classification terms below with the ones that match your work.
%% ACM Digital library provides a taxonomy and a tool for classification
%% in computer science. Use 1-3 paths, and use right arrows between the
%% about three levels in the path; each path requires a new line.

\classification{\protect{\ \\
\  General and reference $\rightarrow$ Document types  $\rightarrow$ Surveys and overviews\  \\
\  Applied computing  $\rightarrow$ Document management and text processing  $\rightarrow$ Document management $\rightarrow$ Text editing
}}

%% If you want to quote someone special. You can comment this line out and there will be nothing on the document.
%\quoting{Bachelor's degrees make pretty good placemats if you get them laminated.}{Jeph Jacques}


%% OPTIONAL STEP: Set up properties and metadata for the pdf file that pdfLaTeX makes.
%% Your name, work title, and keywords are recommended.
\hypersetup{
    unicode=true,           % to show non-Latin characters in Acrobat’s bookmarks
    pdftoolbar=true,        % show Acrobat’s toolbar?
    pdfmenubar=true,        % show Acrobat’s menu?
    pdffitwindow=false,     % window fit to page when opened
    pdfstartview={FitH},    % fits the width of the page to the window
    pdftitle={},            % title
    pdfauthor={},           % author
    pdfsubject={},          % subject of the document
    pdfcreator={},          % creator of the document
    pdfproducer={pdfLaTeX}, % producer of the document
    pdfkeywords={something} {something else}, % list of keywords for
    pdfnewwindow=true,      % links in new window
    colorlinks=true,        % false: boxed links; true: colored links
    linkcolor=black,        % color of internal links
    citecolor=black,        % color of links to bibliography
    filecolor=magenta,      % color of file links
    urlcolor=cyan           % color of external links
}

%%-----------------------------------------------------------------------------------

\begin{document}

% Generate title page.
\maketitle

%%%%%%%%%%%%%%%%%%%%%%%%%%%%%%%%%%%%%%%%%%%%%%%%%%%%%%%%%
%% STEP 3:
%%%%%%%%%%%%%%%%%%%%%%%%%%%%%%%%%%%%%%%%%%%%%%%%%%%%%%%%%
%% Write your abstract in the separate file, to be positioned here.
%% You can make several abstract pages (if you want it in different languages),
%% in which case you should also define the language of the abstract,
%% as below.

% \begin{abstract}{finnish}

% Tämä dokumentti on tarkoitettu Helsingin yliopiston tietojenkäsittelytieteen osaston opin\-näyt\-teiden ja harjoitustöiden ulkoasun ohjeeksi ja mallipohjaksi. Ohje soveltuu kanditutkielmiin, ohjelmistotuotantoprojekteihin, seminaareihin ja maisterintutkielmiin. Tämän ohjeen lisäksi on seurattava niitä ohjeita, jotka opastavat valitsemaan kuhunkin osioon tieteellisesti kiinnostavaa, syvällisesti pohdittua sisältöä.


% Työn aihe luokitellaan  
% ACM Computing Classification System (CCS) mukaisesti, 
% ks.\ \url{https://dl.acm.org/ccs}. 
% Käytä muutamaa termipolkua (1--3), jotka alkavat juuritermistä ja joissa polun tarkentuvat luokat erotetaan toisistaan oikealle osoittavalla nuolella.

% \end{abstract}

\begin{otherlanguage}{english}
\begin{abstract}
Write your abstract here.

In addition, make sure that all the entries in this form are completed.

Finally, specify 1--3 ACM Computing Classification System (CCS) topics, as per \url{https://dl.acm.org/ccs}.
Each topic is specified with one path, as shown in the example below, and elements of the path separated with an arrow.
Emphasis of each element individually can be indicated
by the use of bold face for high importance or italics for intermediate
level.

\end{abstract}
\end{otherlanguage}


% Place ToC
%\newpage
\mytableofcontents

\mainmatter

%%%%%%%%%%%%%%%%%%%%%%%%%%%%%%%%%%%%%%%%%%%%%%%%%%%%%%%%%
%% STEP 4: Write the thesis.
%%%%%%%%%%%%%%%%%%%%%%%%%%%%%%%%%%%%%%%%%%%%%%%%%%%%%%%%%
%% Your actual text starts here. You shouldn't mess with the code above the line except
%% to change the parameters. Removing the abstract and ToC commands will mess up stuff.
%%
%% Command \include{file} includes the file of name file.tex.
%% A new page will be created at every \include command, 
%% which makes it appropriate to use it for large entities such as book chapters. Cannot be nested.
%% It is useful for a big project, as changing one of the include targets 
%% won't force the regeneration of the outputs of all the rest.
%% Alternatively, \input is a more lower level macro 
%% which simply inputs the content of the given file like it was copy&pasted there manually.

\chapter{Introduction\label{intro}}

Within the past decade, web accessibility legislation in the EU has enforced the public sector websites to comply with web accessibility guidelines \citep{eudirective2016}. Additionally, the European Accessibility Act will force some private sector websites to comply after the second half of 2025 \citep{eudirective2019}. The legislation aims to ensure equal access to digital information and services for people with disabilities. However, more than 95\% of the top million websites have accessibility issues that are detectable with Accessibility Evaluation Tools (AETs) \citep{webaimmillions}.

There are multiple AETs to help assess accessibility compliance \citep{govukaccessibility}. However, these tools cover reliably around 20--30\% of the 86 success criteria \citep{govukaccessibilityresults, dequecoverage, webaimmillions}. Therefore, human evaluation and expertise in accessibility are necessary to uncover barriers that affect people with disabilities \citep{wcagevaluationtools}. An accessibility evaluation tool can detect if a website has a title, but a human has to evaluate if the title describes the content of the page. 

This thesis explores web accessibility evaluation and the potential of Large Language Models (LLMs) in web accessibility evaluation, emphasizing 2.4.2 Page Titled Success Criterion that requires human evaluation. A design science research approach is taken to build an input prompt for an LLM to help evaluate if the title describes the content. The thesis is structured as follows.

Chapter 2 introduction to web accessibility, covering definitions, common disabilities affecting web use, implementation, guidelines, and legislation in the EU. Chapter 3 focuses on accessibility evaluation methodologies, accessibility evaluation tools, manual evaluation, and the potential role of AI within accessibility. Evaluation methodologies are presented, following the capabilities of AETs, detailing which criteria they cover, how the effectiveness and reliability of AETs have been studied, and what has been done to standardize the AET's development.

Chapter 4 details the selected research methods, research questions, problem identification, selected tools and scope, and the two iteration phases of the input prompt created in the study, along with the evaluation metrics used. The results of each iteration are presented in Chapter 5, highlighting the accuracy, consistency, and perceived usefulness of the LLM's output.

Chapter 6 summarizes the main findings, discusses the implications, and addresses study limitations. The thesis concludes by emphasizing the importance of web accessibility evaluation through AI-assisted approaches.
\chapter{Web accessibility\label{accessibility}}

Web accessibility aims to ensure an equitable online experience for users of diverse abilities. According to World Health Organization, an estimate of over 1.3 billion people worldwide face some form of significant disability \citep{whodisability}. Persons with low vision, hard of hearing or learning disability might have difficulties perceiving information on the web. The significance of web accessibility is increasingly apparent. This chapter delve into web accessibility. In section 2.1....

\section{Definition}

Web accessibility is defined as an equal opportunity to perceive, operate, understand and interact with an online service \citep{webaccessibilitydefinition}. Regardless of disability a person faces, the same information and functionality should be made available. Web accessibility considers the technical and content implementation. Technical implementation of web accessibility ensures that barriers for assistive technologies, such as screen readers or speech input software, are removed. Structuring content with headings, using short paragraphs and simple understandable language makes the content accessible. Additionally, content accessibility is providing content in multiple forms, such as videos, sound, images and text.

\section{Disabilities affecting the use of web}

Spanning from visual impairments to situational disabilities, each condition presents barriers that hinder the ability to interact and comprehend information on the web. Understanding these disabilities is crucial when designing online services that prioritize inclusion and ensures an equal experience for all users. Disabilities that affect using the web are vision, physical, auditory and cognitive \citep{w3cbarriers}.

One of the most commonly known condition affecting web accessibility is vision impairment and blindness with more then 2.2 billion people worldwide impacted \citep[Chapter~1]{whovision, webaccessibility}. Individuals with visual impairments depend on functionality for formatting the content to their needs. Barriers encountered by individuals with visual impairment are lack of text alternatives for images and controls, insufficient contrast between text color and background color and lack of keyboard navigation support. 

Accessibility is not only targeted for individuals with impairments. Individuals that have lost their glasses, has a broken arm or a hard time reading the screen when the sun glares also benefits of web accessibility. In addition to temporal disabilities and situational limitations, also aging people benefit. Most of the people with low vision are above the age of 50 \citep{whovision}.


\section{Implementation and standards}

Web accessibility strives for an equal opportunity for persons with any sort of impairment to interact with web services on any device. In web accessibility the goal is to implement a web page in a specific way so that assistive technologies can help to remove the barrier for accessing the information provided. Commonly known assistive technologies are screen readers, screen magnifiers and speech input software that are incorporated in major operating systems.


Web designers, developers and content creators implement web accessibility. Accessibility is built upon these three topics, technical implementation, ease of use and comprehensibility of content \citep{webaccessibilitydefinition}. Technical implementation provides functionality for assistive technologies in the web source code based on web standards. Ease of use involves having a user experience with clearly structured frame for navigation and content placement. Comprehensibility of content is the structure of the content itself, meaning the content is written in short paragraphs and plain language. 

In web development the technical implementation is mainly done by web developers by following the HTML and CSS standard, ease of use is web designers responsibility and page content itself is written and structured by content creators. However, there are overlap in responsibilities of implementing accessibility, as creation of web pages has become easier due to website builders and different web frameworks. In the EU, the implementation details for web accessibility have been standardized based on existing guidelines.



\section{Legislation}

The emergence of directives and legislative measures plays a role in the goal to develop best practices for web accessibility. The European Union (EU) Directive 2016/2102 was adopted in 2016 to ratify part of the United Nations Convention on the Rights of Persons with Disabilities (UN CRPD). The Directive enforces EU Member States to ensure that their public sector organizations implement and monitor accessibility on their web pages and mobile applications \citep{eudirective2016}. The directive gave a clear timeline for the member states. At latest of September 2018 each member state should have transposed this directive into national legislation. Since September 2019 all new web sites must conform. All public sector websites had to comply after September 2020. In June 2021 mobile applications had to conform and in December 2021 the monitoring and reporting started. 

The EU Directive adopted in 2016 refers to the European Standard EN 301 549 that specifies the accessibility requirements for ICT products and services. In the first version of the Standard the guidelines and requirements were based on Web Content Accessibility Guidelines provided by the World Wide Web Consortium \citep{wcagadoptioneurope}. During the same time as member states had to transpose this into a national law in 2018, the standard was changed to be a direct reference to the WCAG making the WCAG an accessibility standard that EU member states should follow. There are three levels of web accessibility defined: A being the lowest, AA and AAA being the highest. The directive requires the public sector to conform with AA level of the guidelines.

Monitoring of accessibility in the public sector should be done every three years \citep{eudirectivemonitoring}. The first monitoring and accessibility report from member states to the EU commission was given in 2021. Next reports should be delivered in December 2024. This monitoring is divided into two review categories, in-depth and simplified. Simplified review requires accessibility testing using only automated accessibility evaluation tools, where as in-depth review requires both tools and manual review by accessibility specialists. 

Based on the public sector directive, the EU Commission has created a similar directive affecting some private sector bodies. The European Accessibility Act (EAA) was adopted by the EU in 2019 to complement the 2016 directive and UN CRPD \citep{eudirective2019}. The EAA will require medium and large sized private sector companies in the field of banking, travel and e-commerce starting to follow accessibility guidelines from June 2025 forward.

\section{Guidelines}

The World Wide Web Consortium creates accessibility guidelines through the Web Accessibility Initiative \citep{wcagoverview}. The international standard is named Web Accessibility Content Guidelines (WCAG). Version 2 of WCAG was published in 2008. Version 2.1 came out in 2018 and is the version referenced in the European Standard. Version 2.2 was published, and made a recommendation by W3C in October 2023. The W3C recommends always to use the latest version as they are backwards compatible. The EU commission is in progress of harmonising the EN 301 549 Standard with the latest WCAG version with a planned release of 2025 \citep{etsi_standard}.

The guidelines aim to ensure that digital content is accessible to a broad audience and adaptable to various forms, accommodating diverse sensory, physical, and cognitive abilities of individuals \citep{wcag22}. The WCAG document is written in a technology agnostic way. It is based on four principles: \textit{perceivable, operable, understandable, and robust}. In each principle, there are guidelines and under each guideline are success criteria. There are a total of 13 guidelines and 86 success criteria. In total there are three levels of accessibility; A being the lowest, AA, and AAA being the highest, see table \ref{table:wcag22}. 

Each success criterion is part of one level of accessibility. Version 2.2 remove one and add nine new success criteria. Success criterion 4.1.1 is made obsolete, as assistive technologies do not need to directly parse the HTML code anymore, and problems occurred by this criterion is addressed by other criteria. The new nine added success criteria are: 2.4.11, 2.4.12, 2.4.13, 2.5.7, 2.5.8, 3.2.6, 3.3.7, 3.3.8 and 3.3.9 as shown in table \ref{table:wcag22}. Six of these are required to be implemented to achieve AA-level of accessibility as per legal requirements.

\begin{table}[]
\caption{WCAG 2.2 principles, guidelines and success criteria divided with the three accessibility levels. Adapted from \textcite{wcag22}}
\begin{adjustbox}{width=1\textwidth}
\begin{tabular}{|l|l|l|l|l|}
\hline
\textbf{Principle} &
  \multicolumn{1}{c|}{\textbf{Guideline}} &
  \multicolumn{1}{c|}{\textbf{\begin{tabular}[c]{@{}c@{}}Success Criterion\\ Level A\end{tabular}}} &
  \multicolumn{1}{c|}{\textbf{\begin{tabular}[c]{@{}c@{}}Success Criterion\\ Level AA\end{tabular}}} &
  \multicolumn{1}{c|}{\textbf{\begin{tabular}[c]{@{}c@{}}Success Criterion\\ Level AAA\end{tabular}}} \\ \hline
\multirow{4}{*}{\textbf{Perceivable}} &
  1.1 Text alternatives &
  1.1.1 Non-text content &
   &
   \\ \cline{2-5} 
 &
  \begin{tabular}[c]{@{}l@{}}1.2 Time-based media \\ (Prerecorded, \\ unless stated otherwise)  \end{tabular} &
  \begin{tabular}[c]{@{}l@{}}1.2.1 Audio-only and \\ Video-only\\ 1.2.2 Captions\\ 1.2.3 Audio description \\ or media alternative\end{tabular} &
  \begin{tabular}[c]{@{}l@{}}1.2.4 Captions (Live)\\ 1.2.5 Audio description\end{tabular} &
  \begin{tabular}[c]{@{}l@{}}1.2.6 Sign language\\ 1.2.7 Extended audio description\\ 1.2.8 Media alternative\\ 1.2.9 Audio-only (Live)\end{tabular} \\ \cline{2-5} 
 &
  1.3 Adaptable &
  \begin{tabular}[c]{@{}l@{}}1.3.1 Info and relationships\\ 1.3.2 Meaningful sequence\\ 1.3.3 Sensory characteristics\end{tabular} &
  \begin{tabular}[c]{@{}l@{}}1.3.4 Orientation\\ 1.3.5 Identify input purpose\end{tabular} &
  1.3.6 Identify purpose \\ \cline{2-5} 
 &
  1.4 Distinguishable &
  \begin{tabular}[c]{@{}l@{}}1.4.1 Use of color\\ 1.4.2 Audio control\end{tabular} &
  \begin{tabular}[c]{@{}l@{}}1.4.3 Contrast (Minimum)\\ 1.4.4 Resize text\\ 1.4.5 Images of text\\ 1.4.10 Reflow\\ 1.4.11 Non-text contrast\\ 1.4.12 Text spacing\\ 1.4.13 Content on hover or \\ focus\end{tabular} &
  \begin{tabular}[c]{@{}l@{}}1.4.6 Contrast (Enhanced)\\ 1.4.7 Low or no background audio\\ 1.4.8 Visual presentation\\ 1.4.9 Images of text \\ (No exception)\end{tabular} \\ \hline
\multirow{5}{*}{\textbf{Operable}} &
  \begin{tabular}[c]{@{}l@{}}2.1 Keyboard \\       accessible\end{tabular} &
  \begin{tabular}[c]{@{}l@{}}2.1.1 Keyboard\\ 2.1.2 No keyboard trap\\ 2.1.4 Character key shortcuts\end{tabular} &
   &
  2.1.3 Keyboard (No exception) \\ \cline{2-5} 
 &
  2.2 Enough time &
  \begin{tabular}[c]{@{}l@{}}2.2.1 Timing adjustable\\ 2.2.2 Pause, stop, hide\end{tabular} &
   &
  \begin{tabular}[c]{@{}l@{}}2.2.3 No timing\\ 2.2.4 Interruptions\\ 2.2.5 Re-authenticating\\ 2.2.6 Timeouts\end{tabular} \\ \cline{2-5} 
 &
  \begin{tabular}[c]{@{}l@{}}2.3 Seizures and \\       physical reactions\end{tabular} &
  \begin{tabular}[c]{@{}l@{}}2.3.1 Three flashes or \\          below threshold\end{tabular} &
   &
  \begin{tabular}[c]{@{}l@{}}2.3.2 Three flashes\\ 2.3.3 Animation from interactions\end{tabular} \\ \cline{2-5} 
 &
  2.4 Navigable &
  \begin{tabular}[c]{@{}l@{}}2.4.1 Bypass block\\ 2.4.2 Page titled\\ 2.4.3 Focus order\\ 2.4.4 Link purpose (In context)\end{tabular} &
  \begin{tabular}[c]{@{}l@{}}2.4.5 Multiple ways\\ 2.4.6 Headings and labels\\ 2.4.7 Focus visible\\ 2.4.11 Focus not obscured \\            {[}New{]}\end{tabular} &
  \begin{tabular}[c]{@{}l@{}}2.4.8 Location\\ 2.4.9 Link purpose (Link only)\\ 2.4.10 Section headings\\ 2.4.12 Focus not obscured \\ (Enhanced) {[}New{]}\\ 2.4.13 Focus appearance {[}New{]}\end{tabular} \\ \cline{2-5} 
 &
  2.5 Input modalities &
  \begin{tabular}[c]{@{}l@{}}2.5.1 Pointer gestures\\ 2.5.2 Pointer cancellation\\ 2.5.3 Label in name\\ 2.5.4 Motion actuation\end{tabular} &
  \begin{tabular}[c]{@{}l@{}}2.5.7 Dragging movements \\ {[}New{]}\\ 2.5.8 Target size \\ (Minimum) {[}New{]}\end{tabular} &
  \begin{tabular}[c]{@{}l@{}}2.5.5 Target size (Enhanced)\\ 2.5.6 Concurrent input mechanisms\end{tabular} \\ \hline
\multirow{3}{*}{\textbf{Understandable}} &
  3.1 Readable &
  3.1.1 Language of page &
  3.1.2 Language of parts &
  \begin{tabular}[c]{@{}l@{}}3.1.3 Unusual words\\ 3.1.4 Abbrevations\\ 3.1.5 Reading level\\ 3.1.6 Pronunciation\end{tabular} \\ \cline{2-5} 
 &
  3.2 Predictable &
  \begin{tabular}[c]{@{}l@{}}3.2.1 On focus\\ 3.2.2 On input\\ 3.2.6 Consistent help {[}New{]}\end{tabular} &
  \begin{tabular}[c]{@{}l@{}}3.2.3 Consistent navigation\\ 3.2.4 Consistent identification\end{tabular} &
  3.2.5 Change on request \\ \cline{2-5} 
 &
  3.3 Input Assistance &
  \begin{tabular}[c]{@{}l@{}}3.3.1 Error identification\\ 3.3.2 Labels or instructions\\ 3.3.7 Redundant entry {[}New{]}\end{tabular} &
  \begin{tabular}[c]{@{}l@{}}3.3.3 Error suggestions\\ 3.3.4 Error prevention \\ (Legal, Financial, Data)\\ 3.3.8 Accessible authentication \\ (Minimum) {[}New{]}\end{tabular} &
  \begin{tabular}[c]{@{}l@{}}3.3.5 Help\\ 3.3.6 Error prevention (All)\\ 3.3.9 Accessible authentication \\ (Enhanced) {[}New{]}\end{tabular} \\ \hline
\textbf{Robust} &
  4.1 Compatible &
  \begin{tabular}[c]{@{}l@{}}4.1.1 Parsing \\ (Obsolete and removed)\\ 4.1.2 Name, role, value\end{tabular} &
  4.1.3 Status messages &
   \\ \hline
\end{tabular}
\end{adjustbox}
\label{table:wcag22}
\end{table}

Perceivable means that information and elements can be perceived by all users \citep{wcag22}. Non-textual content should have a textual representation, such as alternative text for images or textual description for prerecorded audio and video (guidelines 1.1 and 1.2). Adaptable describes that content is structured so that computer programs can convey the information to users regardless of device screen size or orientation. Distinguishable main purpose is for users to be able to distinguish between foreground and background information.

Operable principle ensures that the page can be used with different input peripherals, such as keyboard, mouse or touch screen \cite{wcag22}. Keyboard accessible guideline 2.1 defines that same functionality should be available by only using the keyboard. Guideline 2.2 (Enough time) requires time based content, such as notifications, to be available and controllable by the user. Seizure and physical reactions (2.3) asks designers not to create flashing content and limits the flashing of a page to three. Navigable (2.4) provides users information where they are on the page and quick access to relevant content. A known accessibility feature for navigating is the "Skip to content" link that is provided on pages for screen readers to quickly jump to the content block. The last guideline 2.5 is input modalities that guides developers to ensure functionality beyond keyboard input.

Understandable principle is about the accessibility of the content. Readable (3.1) requires web developers to set language for the page. Predictable expects that there are no sudden changes on the page, such as requesting for help button would change position. Input assistance ensures that input elements have labels, the same information would not be asked twice from the user in the same session and that input requirements and errors are displayed correctly. 

Robust requires the page to be created in such a way that it can be interpreted by assistive technologies now and in the future \citep{wcag22}. Compatible is the sole guideline defining that developers should apply roles and names for elements and that status messages can be read by assistive technologies without focusing on the element. Parsing was removed in version 2.2 as it has been addressed by other success criteria in the documents.

The success criteria under the guidelines are written as testable statements \citep{wcag22}. WCAG also offers guidance and best practices for each success criteria by providing informative techniques for web content developers. The techniques are categorized as sufficient and advisory, where advisory techniques go beyond the requirements for the specific criterion. Common failures for conformance are also documented. The techniques describing sufficient conformance are more technical and contains examples on how to implement a success criterion with HTML and CSS techniques. 

There are in total 55 success criteria to cover for conformance with AA level of web accessibility in WCAG 2.2. A comprehensive checklist \citep{wcag_checklist} of tips and techniques to achieve conformance with AA level of accessibility has been developed by the W3C.

% \section{Accessibility in web programming}

% - Talk about what the guidelines means for a web developer perspective; use WCAG 2.2 "How to meet WCAG compliance" as reference

% - 

% - Dive into some of the most basic accessibility, describe a common case, success criterion 1.1.1 for example


\chapter{Web accessibility evaluation\label{accessibility_evaluation}}

As web services undergo rapid development and continuous deployment, a continuous accessibility evaluation is crucial, mandated by both legislative organizations and end-users. The primary goal of assessing web accessibility is to promote digital inclusion by identifying and eliminating barriers, thereby expanding access to a wider audience. Multiple accessibility evaluation tools are available to help web content developers in ensuring inclusivity and conformance to accessibility standards.

\section{Methodologies}

Web accessibility evaluation is a process of assessing and determining to which extent a website is accessible to people with disabilities \citep[Chapter~26.2]{webaccessibility}. The evaluation of web accessibility involves analyzing each web page on the site against established standards and guidelines. The process of web accessibility evaluation can be categorized into three categories: automated testing, manual inspection, and user testing.

Automated testing is done by accessibility evaluation tools that are programmed to automatically parse and evaluate the source code of a web page based on guidelines \citep[Chapter~26.2]{webaccessibility}. Adopting automated testing early ensures that potential accessibility issues are identified and addressed with immediate feedback during the development cycle. However, automated testing tools can only cover those guidelines that are assessable through machine-based analysis.

Manual inspection is required to ensure that a web page conforms to accessibility standards \citep[Chapter~26.2]{webaccessibility}. A conformance review is the most frequently used methodology that involves evaluators checking if the web page meets criteria based on a checklist. \textcite{comparative_accessibility_methods} introduced and compared his Barrier Walkthrough methodology where evaluation is done in a more systematic way taking into consideration the context of the scenario and goal. An example of the scenario and goal is a person with vision impairment, using a specific assistive technology trying to fill in a form. Results show that the Barrier Walkthrough method is better at finding critical problems more accurately, but fails in finding all the accessibility problems \citep{comparative_accessibility_methods}. Nevertheless, a manual inspection conducted by different evaluators may result in different outcomes regardless of the methodology used \citep{accessibility_evaluation_experts, 10.1145/1878803.1878813_testability_expertise}. 

User testing is based on empirical usability testing. User testing is the most reliable accessibility evaluation method involving actual users with disabilities performing tasks on a web page \citep[Chapter~26.2]{webaccessibility}. Reliability is the consistency of the outcome regardless of whom performed the test. However, user testing is slow and costly. In addition, it is complex to set up the testing session and to take into account a diverse range of users \citep{comparative_accessibility_methods}

Multiple evaluation methodologies have been proposed combining the three categories described above \citep[Chapter~26.2.1]{webaccessibility}. It is recommended to always use more than one evaluation tool, as the technical implementation of each tool might differ ending up in different results. Manual inspection and user testing should always be done to ensure conformance. However, there is no consensus on how to combine these methods for a comprehensive accessibility study \citep[Chapter~26.2.1]{webaccessibility}.

\section{Automated tools\label{automated_tools}}

Automated Accessibility Evaluation Tools (AET) should be used to check if a web page conforms to accessibility guidelines. An accessibility evaluation tool parses the HTML and CSS code of a web page and checks that the web page conforms to the accessibility standards. These tools are helpful for quick accessibility evaluation. However, human evaluation is always required to check for full conformance as these tools might produce misleading results \citep{wcagevaluationtools}. Therefore, accessibility evaluation tools should be used as assistive evaluators when evaluating the conformance of a website.

There are dozens of tools to choose from, differing in functionality, coverage, and features \citep{tool_list}. Studies show that the amount of errors found on a page varies largely between tools \citep{comparison_10.1145/3371300.3383346, comparison_10.1145/3607720.3607722, tool_analysis_directive}. Where one tool can find a dozen errors, another tool can find thousands on the same page. \textcite{tool_analysis_directive} reviewed the monitoring reports from member states and did a comparative analysis on 10 free-to-use tools. Results show that almost all tools link found problems to WCAG success criteria. However, tools differed in the amount of criteria they checked and in the final representation of the coverage report. Furthermore, the transparency of tools on which success criteria they cover is not clear to users \citep{tool_analysis_directive}.

A method to compare and measure tool effectiveness was proposed by \textcite{Brajnik2004}. The effectiveness of a tool can be measured with correctness, completeness, and specificity. Correctness measures how tools report non-existing problems (false positives). Completeness measures how tools fail to find problems that exist (false negatives). A tool can be considered highly specific if it can detect a wide range of distinct accessibility issues, offering detailed insights into each problem. \textcite{benchmark_aet} studied the coverage, completeness, and correctness of six popular tools showing that a higher amount of accessibility issues on a page gives a higher completeness score, and in contradiction, on highly accessible pages the completeness score drops. In addition, the study shows that tools with high completeness scores report more easily false positives, reducing the correctness \citep{benchmark_aet}. \textcite{tooltransparency} evaluated four tools on their specificity and transparency. Results show that three of the tools share what success criteria and techniques they cover. However, differences were reported on how tools display the end result to the user.

To mitigate the uncertainty and instability of accessibility evaluation tools the W3C has created a task force working on Accessibility Conformance Testing (ACT) rules \citep{act_overview}. The purpose of ACT rules is to standardize the interpretation of WCAG documents by creating technology-specific test cases that accessibility tool developers and accessibility methodology developers can use to test their outcomes. For unified reporting, the W3C has developed a resource description framework named Evaluation and Report Language (EARL). EARL is used to report conformance to the ACT rules. In addition, it can be used to combine results from different tools. However, currently, there is no standard on how to report the output of an automatic evaluation tool and current tools differ in the reporting style \citep{tool_analysis_directive}.

\section{Tool coverage\label{coverage}}

Automated Accessibility Evaluation Tools (AET's) cover only around 20--30\% of all success criteria in the WCAG \citep{govukaccessibilityresults, dequecoverage, webaimmillions}. In a study conducted by \textcite{govukaccessibilityresults} they created a page with 142 accessibility issues and analyzed this page with 12 different tools (Nu HTML Checker excluded, not an AET). The lowest score was 17 \% by Google Accessibility Developer Tool and the highest scoring tool was SortSite with 40 \%. In addition, the coverage of all tools combined managed to find only 100 out of the 142 accessibility barriers, which supports the need to use multiple tools together when evaluating accessibility \citep{govukaccessibilityresults_blog}. \textcite{comparison_10.1145/3371300.3383346} conducted similar research where they found out that evaluating a page with multiple tools raises the coverage by 10--40 \%. However, \textcite{comparison_10.1145/3371300.3383346} did not study how to combine and remove duplicates from the results between each tool. One large developer and vendor of accessibility testing, Deque Systems, believes that the apprehension on the coverage should be based on real findings rather than which success criteria are covered \citep{dequecoverage}. 

Deque states that up to 57 \% of accessibility issues can be found when considering that there are usually multiple violations for one success criterion on a page \citep{dequecoverage}. From over 2000 in-depth accessibility evaluations they showed that in most cases automation finds more issues than manual review. Data in this study is based on A and AA levels of accessibility when WCAG 2.1 was the recommendation by W3C. The obsolete success criterion 4.1.1 Parsing in WCAG 2.2 was one of the most detected by automation with a proportion of 90.28 \% being found automatically. The six types of issues encountered with a significantly high proportion (percentage in brackets) found by automation are the following: 

\begin{itemize}
  \item 3.1.1 Language of page (91.81 \%, 1 995 issues)
  \item 1.4.3 Contrast (Minimum) (83.11 \%, 73 733 issues)
  \item 2.4.1 Bypass blocks (79 \%, 2 001 issues)
  \item 1.1.1 Non-text context (67.57 \%, 16 014 issues)
  \item 4.1.2 Name, role, value (54.42 \%, 26 276 issues)
  \item 1.3.1 Info and relationship (45.17 \%, 16 432 issues)
\end{itemize}

These six success criteria account for up to 52 \% of all accessibility issues found by automation. By removing the obsolete success criterion 4.1.1 from the data the total amount of found issues by automation is 53 \%. However, this study does not account for the six newly added success criteria in WCAG 2.2 Level AA. Moreover, \textcite{dequeaxe4_5} writes in Deque's blog that the only identified success criterion to be testable in WCAG 2.2 according to their promise, that is not returning false positives, is the success criterion 2.5.8 Target size \citep{dequeaxe4_5}. Deque's target is to ensure that the axe-core engine reports zero false positives to ensure that the results can be trusted by developers and accessibility evaluators \citep{dequecoverage}. Axe-core is an open-source accessibility testing engine for web browsers. It is used by millions of Github projects and it also works as the core for Deque's tools, and Google Lighthouse. 

As WCAG version 2.2 is the new standard, according to the coverage report by \textcite{dequecoverage} there would be 17 out of 86 success criteria that automation can discover with certainty. An automated accessibility evaluator can determine that a page has a title or an image has an alternative text from the web page source code. However, these tools can not determine if a page title or alternative text for non-text content is descriptive \citep{wcag_checklist}.

\section{Manual evaluation}

Automated accessibility evaluation tools are the first step in finding accessibility problems on a web page. To increase the coverage, manual inspection is required. Manual evaluation is a costly and time-consuming process as each web page on the whole site has to be evaluated separately. Most of the success criteria in the WCAG can only be determined properly by human evaluation, such as criteria related to context. For example, test automation can detect if the alternative (alt) attribute is set for image tags in the HTML code. However, it can not determine if the alternative text describes the image correctly to users \citep{comparison_10.1145/3371300.3383346}. Additionally, expertise does matter when evaluating a page for conformance \citep{10.1145/1878803.1878813_testability_expertise, comparative_accessibility_methods}. In comparison to a novice, an expert in the field of web accessibility is more capable of finding accessibility barriers on a web page.

As there is no standard methodology for manual evaluation, each evaluator has their own toolset that they use to evaluate individual pages manually. Using a semi-automated test tool is a great way to increase coverage of found barriers and to guide the evaluator in particular with wizard based steps guiding the manual evaluation. Semi-automated test tools are used to evaluate a single page in a specific state for accessibility barriers. These tools run automated accessibility evaluations combined with, for example, user input wizards to cover more possible manually evaluated accessibility barriers. 

Two large accessibility evaluation consultant companies, Deque and Siteimprove, have their own paid semi-automated accessibility evaluation tools. \textcite{dequecoverage_semi} extended the same in-depth coverage study mentioned in Section \ref{coverage} on their semi-automated accessibility testing tool. They discovered that the coverage of found accessibility barriers is increased by 23\% when using their wizard based semi-automated accessibility evaluation tool \citep{dequecoverage_semi}. The guided testing is incorporated into their test automation to help evaluators with a more systematic conformance evaluation. For example, trying out Deque's tool free version subscription, the guided testing prompts to check if an alternative text for an image or page title is descriptive in a wizard to the evaluator. The page title wizard was the following question: \blockquote{The page title is 'Understanding WCAG 2.2 | WAI | W3C'. Does it accurately describe the purpose of the page?}. This question leaves the decision up to the evaluator to answer. To answer the question reliably, the evaluator needs to read the content of the page and form an understanding of the entire page.

\section{Use of AI}

With the emergence of multiple Generative AI tools for different contexts, such as music, image, or text generation, the use of AI and machine learning models in web accessibility has sparked discussion and gained focus in the field of accessibility. Companies within the accessibility evaluation industry, such as AudioEye and Deque, have incorporated AI features in their tools \citep{deque_igt, boia_improve_accessibility}. In addition, the current investment in AI has started a cautious discussion within the W3C WCAG Working Group where they are following the progress and possible use-cases of AI in web accessibility \citep{ai_wcag_email}. 

One of the most commonly found issues is a missing descriptive alternative text for images \citep{webaimmillions, dequecoverage}. Therefore, the potential of AI in generating the alternative text for images is an interesting topic with multiple viewpoints \citep{ai_wcag_email, boia_alt_text, potential_for_ai}. Mozilla is experimenting with a machine learning model that runs locally on the user's computer to generate missing image alternatives when viewing PDF files within the browser \citep{alt_image_mozilla}. Even though there are tools to generate automatic alternative descriptions to images, the descriptions are not always considered appropriate as they do not take into account context around the image \citep{accessibility_and_ai, boia_alt_text}. Besides, within the accessibility community, some say that AI should work as an assistant that should not make automatic decisions regarding accessibility \citep{ai_wcag_email, accessibility_and_ai}.

A case study conducted by \textcite{10.1145/3594806.3596542_case_study_gpt} showed that the use of ChatGPT in solving accessibility issues was able to provide correction to code and speed up the process of solving accessibility barriers. However, the study was conducted on two websites and it had problems to solve subjective issues, such as the color contrast between the background color and text color. \textcite{Lopez2024Turning} conducted an experiment where they tested if a LLM would be able to evaluate three different WCAG success criteria that require manual evaluation. In total they had 39 ACT cases and a LLM was able to successfully evaluate 34 of them. However, they had to modify the prompt manually for more complex test cases to specify where to look for the correct information in the HTML tags.


\chapter{Research methods\label{methods}}

This thesis will use a qualitative method based on an iterative manner of the Design Science methodology \citep{designsciencemethodology, iterativedesignscience}. The main goal of this research is do a literature review to identify a problem in a specific context and attempt to create an artifact as a solution to found problem. The context of this thesis will be to research web accessibility evaluation tools and to identify how Generative AI could assist accessibility evaluators to resolve barriers on the web. 

<Insert figure here from \cite{iterativedesignscience}>

The research questions for the thesis are following:

\begin{itemize}
    \item \textbf{RQ1.} What are the limitations of web accessibility evaluation tools in assessing compliance with Web Content Accessibility Guidelines (WCAG)?
    \item \textbf{RQ2.} How well does Generative AI assist to address these limitations?
    \item \textbf{RQ3.} How could these features be incorporated in existing accessibility evaluation tools?
\end{itemize}

To identify current limitations for RQ1 a thorough literature review was conducted with focus on web accessibility from a legislative perspective in the EU.  Chapter \ref{accessibility} describes the progress done to ensure an accessible web. Based on findings from chapter \ref{accessibility} the following search string was used in the ACM Digital Library: "accessibility evaluation" OR "evaluation of tools" OR "comparing tools". Filters applied were: Research Article and Publication Date between 2019 - 2024 to find currently relevant publication within the field. A total of 150 research articles were found. Found papers were selected based on their title and abstract. An inclusion criteria was used based on if either title or abstract mentioned comparing tools, the EU Directive or the WCAG guidelines. Exclusion criteria was used if the paper title or abstract had mentioned one of the following: specific disability, specific technology, mobile accessibility or the study was conducted in a country. A snowballing method was used to identify the most commonly referred papers which are also used in this thesis as references. Additionally, found researches were searched for to find relevant literature.






- Literature review chapter 4, get to know legislation from EU perspective, directives and memos in directive, read about guidelines -> EU Standard -> WCAG, explain accessibility and the accessibility guidelines in an understandable manner as short as possible with good literature (EU / State web pages)

- Literature review on accessibility evaluation, "accessibility evaluation tools" in google scholar, use ACM as main source, search "web accessibility" OR "accessibility evaluation" OR "evaluation of tools", filter by relevance and read papers, identify common themes and references, snowball method through recent papers to find the current themes and referred papers

\chapter{Results\label{results}}

This chapter presents the results and answers the research questions. Section \ref{rq1_result} answers the first research question. In Section \ref{rq2_result} the outcome of the LLM is analyzed for each iteration based on the evaluation criteria. In the last section, the artifact is studied from the perspective of how it could be used in other possible success criteria and ACT rules provided by W3C.

\section{Limitations of web accessibility evaluation tools\label{rq1_result}}

The literature review done in Chapter \ref{accessibility_evaluation} answers RQ1. In total, there are 86 success criteria in the WCAG 2.2 documents. Out of these 86 success criteria, test automation covers only 17 reliably (zero false positives). Using manual evaluation with semi-automated web accessibility evaluation tools increases the coverage, and adds to the sufficiency of these 17 found by automation \citep{dequecoverage_semi}.

Research on Automated AETs shows that there is a significant distribution of the amount of found accessibility barriers on the same page \citep{comparison_10.1145/3371300.3383346, tool_analysis_directive}. Studies conducted point out that even experts evaluating the same page can end up with different results \citep{10.1145/1878803.1878813_testability_expertise}. Additionally, the transparency of results by AETs to accessibility evaluators even further adds to the complexity of interpreting which success criteria have been sufficiently evaluated.

Semi-automated web accessibility evaluation tools help in conformance evaluation by guiding the evaluator with wizards to further increase coverage \citep{dequecoverage_semi}. However, the outcome of the final evaluation is subjective to the evaluator \citep{10.1145/1878803.1878813_testability_expertise}. To fully conform with some of the success criteria requires knowledge of the context of the web page. Therefore, expertise matters when evaluating web pages for conformance using semi-automated AETs.

\section{Capabilities of LLM's for Page Titled evaluation\label{rq2_result}}

This section describes the outcomes of the artifact and answers RQ2. Some of the success criteria have context-based evaluation metrics, such as 2.4.2 Page Titled success criterion. The results show that the Large Language Model is capable of assisting in a conformance evaluation when evaluating the context-based success criterion 2.4.2 Page Titled. However, a more detailed look at how it ended up with the correct outcome shows a discrepancy.

\begin{table}
\centering
\caption{Table summarizing the overall result of whether the LLM succeeded in the conformance evaluation.}
\label{table:result_summary}
\begin{tabular}{|l|cc|cc|cc|cc|cc|cc|cl|}
\hline
Test case  & \multicolumn{2}{c|}{P1}    & \multicolumn{2}{c|}{P2}    & \multicolumn{2}{c|}{P3}    & \multicolumn{2}{c|}{F1}    & \multicolumn{2}{c|}{F2}    & \multicolumn{2}{c|}{F3}    & \multicolumn{2}{c|}{I1}    \\ \hline
Iteration  & \multicolumn{1}{c|}{1} & 2 & \multicolumn{1}{c|}{1} & 2 & \multicolumn{1}{c|}{1} & 2 & \multicolumn{1}{c|}{1} & 2 & \multicolumn{1}{c|}{1} & 2 & \multicolumn{1}{c|}{1} & 2 & \multicolumn{1}{c|}{1} & 2 \\ \hline
Sequence 1 & \multicolumn{1}{c|}{\cmark} & \cmark & \multicolumn{1}{c|}{\cmark} & \cmark & \multicolumn{1}{c|}{\cmark} & \cmark & \multicolumn{1}{c|}{}  & \cmark & \multicolumn{1}{c|}{}  & \cmark & \multicolumn{1}{c|}{} & \cmark & \multicolumn{1}{c|}{\cmark} &   \\ \hline
Sequence 2 & \multicolumn{1}{c|}{\cmark} & \cmark & \multicolumn{1}{c|}{\cmark} & \cmark & \multicolumn{1}{c|}{\cmark} & \cmark & \multicolumn{1}{c|}{\cmark} & \cmark & \multicolumn{1}{c|}{\cmark} & \cmark & \multicolumn{1}{c|}{\cmark} & \cmark & \multicolumn{1}{c|}{}  & \cmark \\ \hline
Sequence 3 & \multicolumn{1}{c|}{\cmark} & \cmark & \multicolumn{1}{c|}{\cmark} & \cmark & \multicolumn{1}{c|}{\cmark} & \cmark & \multicolumn{1}{c|}{\cmark} & \cmark & \multicolumn{1}{c|}{\cmark} & \cmark & \multicolumn{1}{c|}{\cmark} & \cmark & \multicolumn{1}{c|}{\cmark} & \cmark \\ \hline
Sequence 4 & \multicolumn{1}{c|}{\cmark} & \cmark & \multicolumn{1}{c|}{\cmark} & \cmark & \multicolumn{1}{c|}{\cmark} & \cmark & \multicolumn{1}{c|}{\cmark} & \cmark & \multicolumn{1}{c|}{}  & \cmark & \multicolumn{1}{c|}{\cmark} & \cmark & \multicolumn{1}{c|}{}  &   \\ \hline
Sequence 5 & \multicolumn{1}{c|}{\cmark} & \cmark & \multicolumn{1}{c|}{\cmark} & \cmark & \multicolumn{1}{c|}{\cmark} & \cmark & \multicolumn{1}{c|}{}  & \cmark & \multicolumn{1}{c|}{}  & \cmark & \multicolumn{1}{c|}{\cmark} & \cmark & \multicolumn{1}{c|}{}  &   \\ \hline
\end{tabular}
\end{table}

In both iterations of the artifact, the passing test cases P1, P2, and P3 were the most consistent, accurate, and useful. The LLM did not suggest any improvements for the ACT passing test cases, as each sequence was correctly evaluated and already had a descriptive title present. 

The LLM had issues evaluating the test cases F1 and F2 in the first iteration. The second iteration outcome showed improvements in the conformance evaluation, see Table \ref{table:result_summary}. However, even though the second iteration explicitly states partial conformance for the F1 and F2 test cases in the outcome, how the LLM ended up to this conclusion had fallacies. In both iterations, out of the 30 sequences, there were only two output sequences (F1-1-1 and F1-1-5) that did not suggest an improved title for the page. The F3 test case did not cause major problems for the LLM.

In both iterations of the artifact, only two inapplicable sequences were evaluated correctly, therefore the evaluation outcome can not be relied upon.

Each iteration is analyzed in its subsection in more detail. Iteration results are evaluated more deeply based on usefulness, accuracy, and consistency. See Appendix \ref{appendix:iterations} for actual outputs.

\subsection{First iteration}

The LLM was capable of evaluating if the title was descriptive based on the content of the web page. However, problems occurred in failing and inapplicable test cases. Passing test cases had a successful outcome based on our evaluation criteria, whereas failing and inapplicable test cases had a variance in consistency and fallacies in reasoning.

\subsubsection{Passed test cases}

When assessing the output in detail for usefulness, the LLM provides reasoning for each four of the rules in the input. However, the P2 test case's fourth sequence (P2-1-4) states that the code snippet has only one title element, even though the test case has two. 

In the P3 test case, it can be noted that the LLM is not capable of identifying correctly where the title element is present, and in three out of five sequences, it falsely states that the title element is within the head section of the HTML document, even though it is in the body section. However, this is a minor mistake, as most browsers can fix and set the title correctly.

The accuracy of the accessibility evaluation is described in a short paragraph at the end of the output for all the test cases. Each test case and its sequences have the same outcome from a contextual perspective, therefore the LLM is accurate and consistent when evaluating these passed test cases.

\subsubsection{Failing test cases}

Inconsistency and inaccuracy are observed in the failing test cases F1 and F2 sequences. Additionally, the usefulness of the outcome is not at an appropriate leave as the reasoning has fallacies. However, the F3 test case sequence outcomes did not have similar problems as in test cases F1 and F2.

\subsubsection{F1}

F1-1-2, F1-1-3 and F1-1-4 have the correct outcome where the LLM states partial conformance. Additionally, those sequences that succeed in evaluating conformance correctly also provide suggestions for improvement. 

The first sequence (F1-1-1) has problems understanding that the whole web page code is provided. In the outcome, there is no mention of clementines being the correct context of the page. Additionally, the same uncertainty of what the LLM is evaluating can be understood from the last paragraph which seems to state conformance:

\blockquote{Overall, based on the provided HTML code and the given web accessibility rules, the web page appears to conform with the criteria. However, further evaluation within the context of the entire website and user experience might be necessary for a comprehensive assessment.}

In the fifth sequence (F1-1-5), the LLM explicitly indicates that the page conforms with the 2.4.2 Page Titled success criterion. In addition, the LLM seems to provide an irrational conclusion for the second rule \textit{Check that the title is relevant to the content of the Web page} with an output of:

\blockquote{The title of the web page is "Apple harvesting season". This title seems relevant to the content of the page, which discusses the harvest timing of clementines, a type of fruit.}

\subsubsection{F2}

The F2 test case with two titles, see Figure \ref{failing_cases}, where the first one is incorrect and the second one is correct causes issues with the conformance evaluation made by the LLM. In regards to accuracy, the LLM states conformance in three out of the five sequences. Therefore, three out of five fail to provide an accurate conformance evaluation as F2 is a test case that has partial conformance.

A closer look at the three not correctly evaluated test sequences indicates that the existence of two titles on the web page code, and the fourth rule pointing out to use the first encountered title causes problems for reasoning, see Figure \ref{first_iteration} for the four rules. The LLM finds the second title on the web page and uses it to evaluate conformance for the first three rules. This ordering of rules causes problems for the LLM, indicating that the scope of the fourth rule is not in a logical order.

Below is an example from F2-1-4 output from the LLM on the fourth rule reasoning, which is similar in all three not correctly evaluated cases (F2-1-1, F2-1-4, and F2-1-5): \blockquote{Since browsers typically recognize only the first title element, and in this case, the relevant title is the second one, it aligns with the rule.}. The LLM evaluated the conformance using the second title even though the rule explicitly said to do the evaluation based on the first title found in the HTML code.

Due to the possible problems described above, the accuracy for this test case is not on a sufficient level, nor the consistency. Therefore, the output that the LLM provided for evaluating conformance has fallacies. However, in each sequence, the LLM points out that the second title would be more relevant.

\subsubsection{F3}

All, except the first sequences (F3-1-1), the answer is consistent and all outcomes accurately state partial conformance and provide suggestions on how to improve the title description. The partial conformance from each outcome is based on the first and fourth rule of the input artifact. In other words, the web page has a title and there is only one title present. 

However, the first sequence ends with \blockquote{Overall, while the web page meets the basic requirement of having a title, it could improve its accessibility by making the title more descriptive and directly relevant to the content, such as "Search results for 'accessibility' at University of Arkham".}, indicating that the test case passed the conformance evaluation, but the title could be improved.

\subsubsection{Inapplicable test case (I1)}

In only two of the sequences, it was able to indicate that the page should not be evaluated as it does not have a proper HTML structure as the title tag is used to describe the SVG element. On other test sequences, it satisfied the success criterion. However, as this is a zero-shot prompt there are no examples of how to determine applicability or inapplicability. Therefore, the outcome of the sequences that try to evaluate this inapplicable test case is of limited use if the accuracy can not be improved.

\subsection{Second iteration}

The improvements done in the second iteration of the artifact had a significant effect on the outcome of failing test cases. With the removal of the fourth rule and using it as a separate instruction in the artifact, the usefulness, accuracy, and consistency increased.

\subsubsection{Passed test cases}

In regards to accuracy, all sequences within the passing test cases were evaluated correctly. Therefore, each test case sequence was consistent regarding accuracy.

Logical reasoning in each outcome on how it evaluated conformance is provided by the LLM. Therefore, the artifact and the outcome are useful for the evaluator. As in the first iteration of the artifact, the LLM still had issues identifying in the P3 test case that the title element is not within the head section. As earlier stated in the first iteration results, this is a minor flaw that does not affect the result when the descriptive title is within the body tag.

\subsubsection{Failing test cases}

In comparison to the first iteration of the artifact, significant improvements in overall accuracy can be identified for the failing test cases. Accuracy and consistency improvements are distinguished from the outputs, as each sequence outcome is perceived as partial conformance. 

\subsubsection{F1}

Each of the five sequences ends up in partial conformance. Partial conformance is based on the first rule passing, that is, the web page has a title. One sequence (F1-2-2) stood out in the results, where there is a difference from the other sequences. The F1-2-2 states that it conforms with the third rule where it should check that the web page can be identified by the title \blockquote{Apple harvesting season}, even though the content is about clementines. However, all except the second sequence (F1-2-2) have similar contextual output, where the first rule is met, but the second and third rule fails, ending up in partial conformance.

\subsubsection{F2}

The expected result is that partial conformance should come out of the second and third conditions to check for, as the first title is not relevant to the page content and the title does not identify the web page. 

The input artifact explicitly states the following: \blockquote{The scope of given web accessibility rules are limited to only checking the first title element in a web page.}. However, within each of the outcomes, the LLM has used the second title \blockquote{Clementine harvesting season} in assessing conformance. However, all of them end up in partial conformance due to either the second or third condition being evaluated based on the first title found. Additionally, it can be observed that the second title is always mentioned as a better option for the page title from the LLM output detail of the third rule.

Overall, in terms of accuracy and consistency, each outcome states partial conformance that is described in the last paragraph of each outcome. However, the LLM neglects the information that it should only be evaluating the first found title.

\subsubsection{F3}

A direct quote from one of the sequences summarizes the results very well: \blockquote{Overall, while the page meets the first criterion by having a title, it falls short of the other two criteria as the title isn't directly relevant to the content and may not effectively identify the page.}. Therefore, consistency and accuracy are spot on. Additionally, in each sequence, the reasoning for the outcome is similar and the LLM suggests either explicitly or implicitly a more descriptive title for the test case.

\subsubsection{Inapplicable test case (I1)}

Only the second and third sequence indicates successful evaluation for inapplicability as there is no title element for the whole web page, rather the title element is for the SVG element. For example, the first check that the web page has a title is evaluated in the following way: \blockquote{The web page contains a <title> element within the <svg> tag. However, the <title> element is intended for providing a title for the SVG graphic for accessibility purposes, not for the entire web page}.

\subsection{Generalization}

In Design Science Research Methodologies, the generalization of the artifact is important. Generalization is evaluated on how well the artifact works in other similar contexts \citep{design_science_eval}. This section evaluates the possibility of using the artifact for other conformance-checking techniques provided in the WCAG documents.

In its current form, as a proof of concept, the possibility of using the artifact for other ACT rules is limited because the input prompt is specifically tailored to test if the web page title is descriptive. The artifact is a proof of concept on how LLM evaluates accessibility using the technique provided by \textcite{g88} and modified accessibility support information from \textcite{act_rule_g88}. 

By incorporating the input prompt in a semi-automated accessibility evaluation tool, placeholders in the input prompt could be utilized, therefore improving the generalization of the artifact. However, the artifact works as "the code" for the LLM, and a small change can have a significant change in the quality of the output, as demonstrated in this study. Therefore, placeholders, choice of words, and small changes in the input prompt require thorough testing and evaluation.

Furthermore, currently, the artifact is not tied to any specific accessibility evaluation tool. The input artifact could be used by any organization and in other Large Language Models, as the input prompt was the artifact iterated upon.

To conclude the generalization, the artifact has the potential to be used in other contexts, as the rules could be easily swapped with other ACT rules for other success criteria, and the custom sentence on picking the first title found could be an additional placeholder.

\chapter{Discussion\label{discussion}}

This chapter contains a discussion of the study, implications, limitations, and potential future research. Section 6.1 is a summary of the research questions. Section 6.2 a discussion on the study result and implications are presented. Section 6.3 presents threats to the validity of the study. Section 6.4 promotes potential future research.

\section{Summary of main findings}

Below is a recap of the research questions and a summary of the answers.

\begin{itemize}
    \item \textbf{RQ1.} What are the limitations of web accessibility evaluation tools in assessing compliance with Web Content Accessibility Guidelines (WCAG)?
\end{itemize}

    A multivocal literature review answered the current limitations of web accessibility evaluation tools. Test automation covers 17 out of the 86 success criteria in the WCAG 2.2. However, test automation tools are not capable of thorough evaluation that would meet the WCAG set standard for some of these 17 success criteria. 

    Sufficiently evaluating conformance requires an accessibility specialist evaluation which is a tedious process, as web pages are more complicated than ever. Semi-automated accessibility evaluation tools help evaluators by guiding them through the most common accessibility barriers found on web pages. However, manual evaluation can end up with a different outcome depending on the evaluator's knowledge and workmanship.
    
\begin{itemize}
    \item \textbf{RQ2.} How does Generative AI assist in addressing these limitations?
\end{itemize}

    Generative AI can be utilized to address these limitations. This study shows that Large Language Models (LLM) can assist evaluators in conformance checks on the page title that require a thorough understanding of the content. However, a more detailed observation shows that the LLM had fallacies in how it ended up with the correct outcome when conducting a conformance evaluation for the failing ACT cases. On pages with more content, the outcome of the LLM could speed up the conformance evaluation process. However, as results show, with a zero-shot chain of thought prompting, LLMs are not capable of reliably determining inapplicability.

\section{Study result analysis}

This study shows that even though legislation is moving forward in regards to accessibility, the nature of accessibility and accessibility evaluation is complex, and requires expertise from designers, developers, quality assurance, and accessibility reviewers. An accessibility specialist needs to have a thorough knowledge of the WCAG documents, as no automated or semi-automated accessibility evaluation tool has 100\% coverage. In addition, web developers, designers, and content creators ought to study the same WCAG documents.

An accessibility evaluation tool developer has to understand the ACT rules and sufficient techniques used to check for conformance, as well as how browsers work, to develop a reliable and robust tool for accessibility evaluators. Transparency of evaluation tools helps the evaluator understand which success criteria it covers and to what extent. The study shows that Large Language Models can be utilized to improve the sufficiency of the accessibility evaluation and help ease the conformance evaluation process when using semi-automated accessibility evaluation tools. However, the output of the LLM should be evaluated by a human on how it ended up to the conclusion. Therefore, an LLM would be an additional asset in AETs when conducting conformance evaluations.

\subsection{Prompt iteration}

These results build on existing evidence that LLMs are good zero-shot reasoners \citep{kojima2023large}. With rigorous prompt iteration, the accuracy and quality of the outcome improved. By evaluating the outcome of the LLM, patterns can be detected where the LLM fails to provide reasoning for checks it performed, giving possible directions for improvements. An example observed in this study between the iterations is the order of the conditional checks the LLM should take into account. Therefore, an imperative approach to how the LLM should operate step by step combined with the zero-shot chain of thought improved the quality of the outcome. However, as the second iteration of the artifact went through multiple changes, it remains unknown if \textcite{kojima2023large} zero-shot chain of thought change, or moving the fourth rule as a separate sentence, affected the accuracy and consistency in this study.

\section{Limitations}

The lack of a thorough evaluation, such as surveys or interviews with potential users of this artifact, is a concern of the validity of this study. The evaluation of usefulness is solely based on the observations of the thesis writer. The artifact is a proof of concept and has not been evaluated by accessibility evaluators for usefulness in accessibility conformance reviews. Therefore, a proof of suitability, evaluating whether the LLM would assist and speed up conformance reviews, is yet to be evaluated that would support the findings of the study. 

Two concerns regarding the chosen LLM are that this study was solely done using the ChatGPT 3.5 user interface due to it being free to use and that the LLM is a closed-source tool. Therefore, between the iterations, there is no knowledge if there have been any improvements made by OpenAI to the language model. However, according to OpenAI, the ChatGPT 3.5 model was trained only with data available in early 2022 \citep{openai_35}. Therefore, it can be assumed that the test cases available at \textcite{act_rule_g88} should not be part of the LLM knowledge base. The replication of this study should be straightforward with different LLMs, such as the new GPT-4o, which was announced in September 2024 and is now the main model provided to free-tire users of OpenAI.

Additionally, the characteristics of LLMs, such as non-deterministic output given the same input, or the limited amount of characters that you can input, are a threat to the external validity of the study. This study was conducted using very short code snippets, therefore no input limits were hit. Even though, in the second iteration all the passing and failing test case sequences were correctly evaluated by the LLM, this does not guarantee that the LLM would always correctly evaluate due to the characteristics of LLM being non-deterministic. However, a better design could reduce the input limitation threat. For example, in this study setting, parsing the HTML code, picking the first title element encountered, and the content within the HTML body or main tag would significantly reduce the number of characters sent to the LLM. In addition, the non-deterministic behavior would cause problems if your accessibility evaluation tool promises zero false positives when evaluating for accessibility.

It is beyond the scope of this study to evaluate how the LLM would work if the website language were other than English, or with other large language models available, either provided by some entity online or running the models locally.

\section{Future research}

Future studies should take into account the language used on the website, as accessibility barriers are not only limited to English, specifically, to low-resource languages that have less data available to train the LLMs. In addition, as this study was limited to the short ACT cases, a case study where the artifact would be implemented into a semi-automatic accessibility evaluation tool with a feedback loop from the evaluator would provide insight into how the LLM evaluates more complex websites with more content, and how LLMs could improve the efficiency of the conformance review. Additionally, as conformance reviews are a tedious process, automated ways to measure and observe the output of an LLM's effectiveness and trustworthiness for success criteria evaluations would grow the coverage of automated accessibility evaluation tools.


\chapter{Conclusions\label{conclusions}}

The goal of this study was to figure out how many of the WCAG 2.2 success criteria current Accessibility Evaluation Tools test sufficiently and to see if Large Language Models could complement and be utilized in these AET's. The study shows that with the prompt artifact has an affect on the outcome on the accuracy and consistency of the provided output by the LLM. 




%%%%%%%%%%%%%%%%%%%%%%%%%%%%%%%%%%%%%%%%%%%%%%%%%%%%%%%%%
%\cleardoublepage                          %fixes the position of bibliography in bookmarks
%\phantomsection
\addcontentsline{toc}{chapter}{\bibname}  % This lines adds the bibliography to the ToC
\printbibliography

%%%%%%%%%%%%%%%%%%%%%%%%%%%%%%%%%%%%%%%%%%%%%%%%%%%%%%%%%
\backmatter
\begin{appendices}

%% A sample Appendix
%% 
\appendix{Test case outputs\label{appendix:iterations}}
\addtocontents{toc}{\protect\setcounter{tocdepth}{1}}
The conversation is renamed in form of P1-1-1, where P1 refers to the test case in the \textcite{act_rule_g88}. First number describes the iteration, that is, which version of the artifact and the last number the sequence number ran for the iteration.


\section{P1\label{P1-1}}

\subsection{Iteration 1}

\begin{enumerate}
    \item \href{https://chat.openai.com/share/5589d0f3-a3e3-4aee-bd09-63fb87503f4e}{https://chat.openai.com/share/5589d0f3-a3e3-4aee-bd09-63fb87503f4e}
    \item \href{https://chat.openai.com/share/c9cbb56f-9da3-42f4-928c-7931e1f6d6c6}{https://chat.openai.com/share/c9cbb56f-9da3-42f4-928c-7931e1f6d6c6}
    \item \href{https://chat.openai.com/share/73173b78-220d-49de-85fd-a659f38ce09d}{https://chat.openai.com/share/73173b78-220d-49de-85fd-a659f38ce09d}
    \item \href{https://chat.openai.com/share/b262027f-b852-4ffb-8691-bbff11a7b660}{https://chat.openai.com/share/b262027f-b852-4ffb-8691-bbff11a7b660}
    \item \href{https://chat.openai.com/share/e0c1b60f-b309-4996-a17b-21a3020ee029}{https://chat.openai.com/share/e0c1b60f-b309-4996-a17b-21a3020ee029}
\end{enumerate}

\subsection{Iteration 2}

\begin{enumerate}
    \item \href{https://chatgpt.com/share/10e0e7c2-924f-4c94-b5d3-50fa3ce8e990}{https://chatgpt.com/share/10e0e7c2-924f-4c94-b5d3-50fa3ce8e990}
    \item \href{https://chatgpt.com/share/22b4c9da-ad47-4bed-869b-e36e97d51faf}{https://chatgpt.com/share/22b4c9da-ad47-4bed-869b-e36e97d51faf}
    \item \href{https://chatgpt.com/share/1ffb4ae6-c998-4d38-b935-ac5b603ab482}{https://chatgpt.com/share/1ffb4ae6-c998-4d38-b935-ac5b603ab482}
    \item \href{https://chatgpt.com/share/c42f634e-80a1-44d5-84c2-6ad533804760}{https://chatgpt.com/share/c42f634e-80a1-44d5-84c2-6ad533804760}
    \item \href{https://chatgpt.com/share/d70a12a0-a3c4-410d-86ee-859128972a83}{https://chatgpt.com/share/d70a12a0-a3c4-410d-86ee-859128972a83}
\end{enumerate}

\section{P2\label{P2-1}}

\subsection{Iteration 1}

\begin{enumerate}
    \item \href{https://chat.openai.com/share/fa9ad3e2-ad93-4429-b6ee-38a7ef6bad0e}{https://chat.openai.com/share/fa9ad3e2-ad93-4429-b6ee-38a7ef6bad0e}
    \item \href{https://chat.openai.com/share/79e7dd08-a79e-4594-80ef-87b7a07b64b2}{https://chat.openai.com/share/79e7dd08-a79e-4594-80ef-87b7a07b64b2}
    \item \href{https://chat.openai.com/share/aa354068-bc22-4fdf-9ac3-4a783be9b5d6}{https://chat.openai.com/share/aa354068-bc22-4fdf-9ac3-4a783be9b5d6}
    \item \href{https://chat.openai.com/share/fb200c6f-a83d-4363-81e0-3e869ed52f48}{https://chat.openai.com/share/fb200c6f-a83d-4363-81e0-3e869ed52f48}
    \item \href{https://chat.openai.com/share/af4efc55-de9b-4b0b-b742-074798a1ba0c}{https://chat.openai.com/share/af4efc55-de9b-4b0b-b742-074798a1ba0c}
\end{enumerate}

\subsection{Iteration 2}

\begin{enumerate}
    \item \href{https://chatgpt.com/share/d180c25d-b1da-41eb-8b22-77f1e1b22ac6}{https://chatgpt.com/share/d180c25d-b1da-41eb-8b22-77f1e1b22ac6}
    \item \href{https://chatgpt.com/share/1ef1453a-0349-4aeb-8204-34b65e68a24f}{https://chatgpt.com/share/1ef1453a-0349-4aeb-8204-34b65e68a24f}
    \item \href{https://chatgpt.com/share/c707d684-a4cd-4ff0-94b2-6189302ee5db}{https://chatgpt.com/share/c707d684-a4cd-4ff0-94b2-6189302ee5db}
    \item \href{https://chatgpt.com/share/8e93bbe0-c85e-4a1e-9c3b-5f7aa944e375}{https://chatgpt.com/share/8e93bbe0-c85e-4a1e-9c3b-5f7aa944e375}
    \item \href{https://chatgpt.com/share/9ecfab41-ae15-4ca8-bda3-fdfc969d1842}{https://chatgpt.com/share/9ecfab41-ae15-4ca8-bda3-fdfc969d1842}
\end{enumerate}

\section{P3\label{P3-1}}

\subsection{Iteration 1}

\begin{enumerate}
    \item \href{https://chat.openai.com/share/d2af5c6b-9489-4836-998f-c8af4d99213c}{https://chat.openai.com/share/d2af5c6b-9489-4836-998f-c8af4d99213c}
    \item \href{https://chat.openai.com/share/a7e9bb5a-20ea-4eac-820c-dede05d34279}{https://chat.openai.com/share/a7e9bb5a-20ea-4eac-820c-dede05d34279}
    \item \href{https://chat.openai.com/share/18227ab9-65a6-4eaa-868e-b50774e96199}{https://chat.openai.com/share/18227ab9-65a6-4eaa-868e-b50774e96199}
    \item \href{https://chat.openai.com/share/dabac1ef-8ca3-4765-9366-e9296f9ddd7d}{https://chat.openai.com/share/dabac1ef-8ca3-4765-9366-e9296f9ddd7d}
    \item \href{https://chat.openai.com/share/9dac3e30-e2f0-4ae5-9f9a-20e3d7775d6c}{https://chat.openai.com/share/9dac3e30-e2f0-4ae5-9f9a-20e3d7775d6c}
\end{enumerate}

\subsection{Iteration 2}

\begin{enumerate}
    \item \href{https://chatgpt.com/share/4b52f5f3-2da9-425a-91b7-aef5713bd04f}{https://chatgpt.com/share/4b52f5f3-2da9-425a-91b7-aef5713bd04f}
    \item \href{https://chatgpt.com/share/1c3957f0-aa3c-47cd-9a1e-11d12436b25c}{https://chatgpt.com/share/1c3957f0-aa3c-47cd-9a1e-11d12436b25c}
    \item \href{https://chatgpt.com/share/6731c4b8-9d85-4905-a666-1a8810aafc23}{https://chatgpt.com/share/6731c4b8-9d85-4905-a666-1a8810aafc23}
    \item \href{https://chatgpt.com/share/be4a8fbc-82cc-4408-8507-67187a1b99cf}{https://chatgpt.com/share/be4a8fbc-82cc-4408-8507-67187a1b99cf}
    \item \href{https://chatgpt.com/share/164a9db4-14d8-4cff-a751-a0729b4d4c67}{https://chatgpt.com/share/164a9db4-14d8-4cff-a751-a0729b4d4c67}
\end{enumerate}

\section{F1\label{F1-1}}

\subsection{Iteration 1}

\begin{enumerate}
    \item \href{https://chat.openai.com/share/87a9d162-b402-40fc-8bfa-664c6a5ba3a2}{https://chat.openai.com/share/87a9d162-b402-40fc-8bfa-664c6a5ba3a2}
    \item \href{https://chat.openai.com/share/34809a27-8665-4f6c-84ab-7eb5b243f065}{https://chat.openai.com/share/34809a27-8665-4f6c-84ab-7eb5b243f065}
    \item \href{https://chat.openai.com/share/ccb471b1-522e-4b85-9fec-b405e823086b}{https://chat.openai.com/share/ccb471b1-522e-4b85-9fec-b405e823086b}
    \item \href{https://chat.openai.com/share/af1997ef-ef04-46fc-8cfe-a1f8bba6daab}{https://chat.openai.com/share/af1997ef-ef04-46fc-8cfe-a1f8bba6daab}
    \item \href{https://chat.openai.com/share/c571e9ed-b9ab-4e54-9389-67c456776708}{https://chat.openai.com/share/c571e9ed-b9ab-4e54-9389-67c456776708}
\end{enumerate}

\subsection{Iteration 2}

\begin{enumerate}
    \item \href{https://chatgpt.com/share/07592414-6d62-4929-a9c5-9e51b100f92d}{https://chatgpt.com/share/07592414-6d62-4929-a9c5-9e51b100f92d}
    \item \href{https://chatgpt.com/share/0bbd3995-de1c-407e-a3d2-2f5c730faba3}{https://chatgpt.com/share/0bbd3995-de1c-407e-a3d2-2f5c730faba3}
    \item \href{https://chatgpt.com/share/5d414e15-ea1f-4d5d-900a-612c16f824c1}{https://chatgpt.com/share/5d414e15-ea1f-4d5d-900a-612c16f824c1}
    \item \href{https://chatgpt.com/share/9b7eea10-1ae3-4904-a539-54781329d218}{https://chatgpt.com/share/9b7eea10-1ae3-4904-a539-54781329d218}
    \item \href{https://chatgpt.com/share/b042e503-cd44-488f-ab45-619d63098c1e}{https://chatgpt.com/share/b042e503-cd44-488f-ab45-619d63098c1e}
\end{enumerate}

\section{F2\label{F2-1}}

\subsection{Iteration 1}

\begin{enumerate}
    \item \href{https://chat.openai.com/share/866af13f-2653-4d60-99e5-234f67a2be93}{https://chat.openai.com/share/866af13f-2653-4d60-99e5-234f67a2be93}
    \item \href{https://chat.openai.com/share/1da36f2b-f112-4881-89ad-af37c4ee6aaf}{https://chat.openai.com/share/1da36f2b-f112-4881-89ad-af37c4ee6aaf}
    \item \href{https://chat.openai.com/share/e33b7b58-a43c-4c0b-be08-4e05d31af6f2}{https://chat.openai.com/share/e33b7b58-a43c-4c0b-be08-4e05d31af6f2}
    \item \href{https://chat.openai.com/share/acec840d-8e8b-4a84-8499-4157e62efc41}{https://chat.openai.com/share/acec840d-8e8b-4a84-8499-4157e62efc41}
    \item \href{https://chat.openai.com/share/4bc77120-396e-43a3-8292-8594259b331b}{https://chat.openai.com/share/4bc77120-396e-43a3-8292-8594259b331b}
\end{enumerate}

\subsection{Iteration 2}

\begin{enumerate}
    \item \href{https://chatgpt.com/share/e3e048af-b623-439d-b267-0710f9b0c049}{https://chatgpt.com/share/e3e048af-b623-439d-b267-0710f9b0c049}
    \item \href{https://chatgpt.com/share/fd8d7aab-6447-45f6-aef5-c90fd9171404}{https://chatgpt.com/share/fd8d7aab-6447-45f6-aef5-c90fd9171404}
    \item \href{https://chatgpt.com/share/d5fd65a8-da11-44e5-a065-42ca9455602d}{https://chatgpt.com/share/d5fd65a8-da11-44e5-a065-42ca9455602d}
    \item \href{https://chatgpt.com/share/f05d4764-8b8a-41ab-82b6-a7d83b6bcc85}{https://chatgpt.com/share/f05d4764-8b8a-41ab-82b6-a7d83b6bcc85}
    \item \href{https://chatgpt.com/share/c60d397a-d994-401d-9dc7-89bf3ff00118}{https://chatgpt.com/share/c60d397a-d994-401d-9dc7-89bf3ff00118}
\end{enumerate}

\section{F3\label{F3-1}}

\subsection{Iteration 1}

\begin{enumerate}
    \item \href{https://chat.openai.com/share/86f2fb56-fd98-4791-a764-0d61b13c920f}{https://chat.openai.com/share/86f2fb56-fd98-4791-a764-0d61b13c920f}
    \item \href{https://chat.openai.com/share/f1814074-d83d-429e-ae03-1c6843d4aac0}{https://chat.openai.com/share/f1814074-d83d-429e-ae03-1c6843d4aac0}
    \item \href{https://chat.openai.com/share/08684ef8-34e2-4b72-9d00-5215b0d90e69}{https://chat.openai.com/share/08684ef8-34e2-4b72-9d00-5215b0d90e69}
    \item \href{https://chat.openai.com/share/8cb6a180-8ff0-48e2-aa3c-d5f73f3993e6}{https://chat.openai.com/share/8cb6a180-8ff0-48e2-aa3c-d5f73f3993e6}
    \item \href{https://chat.openai.com/share/7e2f1b2a-a6c7-45f4-9901-7f3cba63c24e}{https://chat.openai.com/share/7e2f1b2a-a6c7-45f4-9901-7f3cba63c24e}
\end{enumerate}

\subsection{Iteration 2}

\begin{enumerate}
    \item \href{https://chatgpt.com/share/2c460c56-930a-45c6-9a9b-9f7e4823fd27}{https://chatgpt.com/share/2c460c56-930a-45c6-9a9b-9f7e4823fd27}
    \item \href{https://chatgpt.com/share/423833b1-db6b-462d-b111-d6872495a113}{https://chatgpt.com/share/423833b1-db6b-462d-b111-d6872495a113}
    \item \href{https://chatgpt.com/share/a3f0e2de-78c6-4205-8275-203514ef4717}{https://chatgpt.com/share/a3f0e2de-78c6-4205-8275-203514ef4717}
    \item \href{https://chatgpt.com/share/ae1d8b0a-738e-491f-921b-5168d8e73bdc}{https://chatgpt.com/share/ae1d8b0a-738e-491f-921b-5168d8e73bdc}
    \item \href{https://chatgpt.com/share/4a04107e-5da9-435d-a61b-416798c929b5}{https://chatgpt.com/share/4a04107e-5da9-435d-a61b-416798c929b5}
\end{enumerate}

\section{I1\label{I1-1}}

\subsection{Iteration 1}

\begin{enumerate}
    \item \href{https://chat.openai.com/share/d6c95007-b965-4ca9-b176-4e3f9899bb0d}{https://chat.openai.com/share/d6c95007-b965-4ca9-b176-4e3f9899bb0d}
    \item \href{https://chat.openai.com/share/a2789f63-0ad6-4e56-9a3f-0fc2de56fdae}{https://chat.openai.com/share/a2789f63-0ad6-4e56-9a3f-0fc2de56fdae}
    \item \href{https://chat.openai.com/share/63d5dc58-6c52-4055-9357-b6e97fbbeeb0}{https://chat.openai.com/share/63d5dc58-6c52-4055-9357-b6e97fbbeeb0}
    \item \href{https://chat.openai.com/share/20b42730-76de-4345-a6db-ef130d43b1a9}{https://chat.openai.com/share/20b42730-76de-4345-a6db-ef130d43b1a9}
    \item \href{https://chat.openai.com/share/9ddfd254-5917-430f-8984-8610e6608c70}{https://chat.openai.com/share/9ddfd254-5917-430f-8984-8610e6608c70}
\end{enumerate}

\subsection{Iteration 2}

\begin{enumerate}
    \item \href{https://chatgpt.com/share/c0148e6e-1121-4f12-94b1-d2a0ae89a37c}{https://chatgpt.com/share/c0148e6e-1121-4f12-94b1-d2a0ae89a37c}
    \item \href{https://chatgpt.com/share/6c7f11fa-59b5-4ed7-915a-379f22cca3a2}{https://chatgpt.com/share/6c7f11fa-59b5-4ed7-915a-379f22cca3a2}
    \item \href{https://chatgpt.com/share/71dfe5ee-5af6-4729-8fe5-09b436f76fb6}{https://chatgpt.com/share/71dfe5ee-5af6-4729-8fe5-09b436f76fb6}
    \item \href{https://chatgpt.com/share/d2f6f76b-47b6-4937-bdb3-a626f9c7b84c}{https://chatgpt.com/share/d2f6f76b-47b6-4937-bdb3-a626f9c7b84c}
    \item \href{https://chatgpt.com/share/312537cb-477d-47f1-8701-dad5a0fbe6ce}{https://chatgpt.com/share/312537cb-477d-47f1-8701-dad5a0fbe6ce}
\end{enumerate}
%% another appendix
%% 
\appendix{Instructions for LaTex}

\section{General Setup}

In the HY-CS-main.tex file you will find the following STEPS 0--5. Below you can find related instructions.
\vspace{0.5cm}

\textbf{STEP 0 -- Access the thesis template}

\begin{itemize}
\item Import the thesis template into a new Overleaf project. The easiest way to do it is to:
\begin{itemize}
    \item Obtain a zip file of the LaTeX template from the webpage of your programme.
    \item Go to \url{https://www.overleaf.com/edu/helsinki} and login to Overleaf with your university credentials.
    \item Go to the list of your projects at \url{https://www.overleaf.com/project}, click ``New Project'' and ``Upload Project''.,  the projects under your account 
    \item Then upload the zip with the template.
    \item You are now ready to write your thesis in Overleaf by editing the template, you can start by renaming the project.
\end{itemize}
\end{itemize}


{\textbf{STEP 1 -- BSc or MSc thesis?}}
\begin{enumerate}
\item Select whether your are writing BSc (tkt) or MSc (csm for CS) thesis.
\item Select your language: \texttt{finnish}, \texttt{english}, or \texttt{swedish}.
\item If you are writing MSc select your line / track.
\end{enumerate}


{\textbf{STEP 2 -- Set up your personal information}}

\begin{enumerate}
\item Specify the title of your thesis with \texttt{\textbackslash title\{\}}.
\item Specify your name to the author field with \texttt{\textbackslash author\{\}}.
\item Specify the names of your supervisors of the thesis with \texttt{\textbackslash supervisors\{\}}.
\item Specify the keywords of the thesis with \texttt{\textbackslash keywords\{\}}.
\item Specify the ACM classification terms of the thesis with \texttt{\textbackslash classification\{\}}. See \url{https://dl.acm.org/ccs} for more information.
\end{enumerate}

{\textbf{STEP 3 -- Write your abstract}}

\begin{itemize}
\item You can have the abstract in multiple languages with the \texttt{otherlanguages} environment. The example below shows how to provide an English abstract: 

\begin{verbatim}
\begin{otherlanguage}{english} 
\begin{abstract}
Your abstract text goes here. 
\end{abstract} 
\end{otherlanguage}
\end{verbatim}

\end{itemize}

{\textbf{STEP 4 -- Writing your thesis}}

\begin{enumerate}
\item There are some minimal contents and instructions below 
\item Remove, or comment out, this appendix from your thesis.
\end{enumerate}

{\textbf{STEP 5 -- Set your bibliography style}}

\begin{itemize}
\item The default is Author-Year style (Einstein, 1905), but it can be easily changed to numbered [1] or alphabetical [Ein05] , as the examples of these are in comments.
\item Discuss the style to use with your supervisor.
\end{itemize}

\section{Bibliography in Latex}

The bibliography is defined in a separate \texttt{.bib} file. For this template, it is named \texttt{bibliography.bib} and includes the content show in Figure~\ref{bibexamples}.

Chapter Bibliography lists all the works that you refer to in your text. You refer to the works in the bibliography using an appropriate \emph{citation key}.
%
%This thesis template contains an example of a bibliography.


References are done using \texttt{\textbackslash citep\{einstein\}}, which generates in text a citation formatted according to the selected style \citep{einstein}, or \texttt{\textbackslash citep\{latexcompanion,knuth99\}}, which generates \citep{latexcompanion,knuth99}. 
As examples of a different kinds of citations (see how these look in the Latex source), we can write \citep{einstein} to refer to the work written by \citeauthor{einstein} in \citeyear{einstein}, because the work by \citet{einstein} appears in the bilbliography included in this template.

Note that there are different possible styles for the bibliography and citation keys.
%
Consult your supervisors on the chosen style -- and once you arrive at a preferred style, use it consistently throughout the thesis.

\begin{figure}[ht]
    \centering
    \begin{scriptsize}
\begin{verbatim}
@article{einstein,
    author =       "Albert Einstein",
    title =        "{Zur Elektrodynamik bewegter K{\"o}rper}. ({German})
        [{On} the electrodynamics of moving bodies]",
    journal =      "Annalen der Physik",
    volume =       "322",
    number =       "10",
    pages =        "891--921",
    year =         "1905",
    DOI =          "http://dx.doi.org/10.1002/andp.19053221004"
}
 
@book{latexcompanion,
    author    = "Michel Goossens and Frank Mittelbach and Alexander Samarin",
    title     = "The \LaTeX\ Companion",
    year      = "1993",
    publisher = "Addison-Wesley",
    address   = "Reading, Massachusetts"
}

@book{knuth99,
    author    = "Donald E. Knuth",
    title     = "Digital Typography",
    year      = "1999",
    publisher = "The Center for the Study of Language and Information",
    series    = "CLSI Lecture Notes (78)"
}\end{verbatim}
\end{scriptsize}
    \caption{Examples of bibliographic reference in .bib file.}
    \label{bibexamples}
\end{figure}

%In the last reference url field the code \verb+%7E+ will translate into \verb+~+ once clicked in the final pdf.

\section{Some instructions about writing in Latex}

The following gives some superficial instructions for using this template for a Master's thesis. For guidelines on thesis writing you can consult various sources, such as university courses on scientific writing or your supervisors.

For more detailed instructions, just google, e.g., "Overleaf table positioning", and your chances of finding good info are pretty good.  


\section{Figures}
Besides text, here are simple examples how you can add figures and tables in your thesis.
Remember always to refer to each figure in the main text and provide them with a descriptive caption.

Figure~\ref{fig:logo} is an example of a figure in the document (see the source about how to add them). 
%Using figures is particularly useful to display plots of experimental results.

\begin{figure}[ht] 
\begin{center}
\includegraphics[width=0.3\textwidth]{template/figures/HY-logo-ml.png}
\caption{University of Helsinki flame-logo for Faculty of Science.\label{fig:logo}}
\end{center}
\end{figure}

\section{Tables}

Table~\ref{table:results} gives an example of a table.
Remember always to cite the table in the main text, table captions go on top of the table. 

\begin{table}[h] % h positions the table here, t! would force on top of the page, or example.
\begin{center}
\caption{Experimental results.\label{table:results}} % caption is here to make it on top
\begin{tabular}{l||l c r} 
Experiment & 1 & 2 & 3 \\ 
\hline \hline 
$A$ & 2.5 & 4.7 & -11 \\
$B$ & 8.0 & -3.7 & 12.6 \\
$A+B$ & 10.5 & 1.0 & 1.6 \\
\hline
%
\end{tabular}
\end{center}
\end{table}



%% yet another appendix
%% 
\appendix{Tutkielmapohjan käyttöohjeet}
\label{appendix:instructions_finnish}

\section{Ensiaslkeleet}

\texttt{HY-CS-main.tex} tiedosto sisältää viisi askelta STEPS 0--5. Alla on kuvattu, mitä nämä askeleet tarkoittavat ja miten niitä seuraamalla luot pohjan tutkielmallesi.
\vspace{0.5cm}

\textbf{STEP 0 -- Kopioi tutkielmapohja}

\begin{itemize}
\item Hae tutkielmapohja uuteen Overleaf-projektiin. Tämä käy helpoiten seuraavasti:
\begin{itemize}
    \item Lataa Latex-pohjan zip-tiedosto koulutusohjelman sivuilta.
    \item Mene osoitteeseen \url{www.overleaf.com/edu/helsinki} ja kirjaudu Overleafiin yli\-opiston tunnuksillasi.
    \item Overleafissa (\url{https://www.overleaf.com/project}), klikkaa ``New Project'' and ``Upload Project''.
    \item Valitse lataamasi tutkielmapohjan zip-tiedosto.
    \item Nyt voit lähteä kirjoittamaan tutkielmaasi suoraan pohjaan, voit aloittaa esim. vaihtamalla projektin nimen.
\end{itemize}
\end{itemize}


{\textbf{STEP 1 -- BSc vai MSc tutkielma?}}
\begin{enumerate}
\item Valitse (tiedostossa \texttt{HY-CS-main.tex}) oletko tekemässä BSc (tkt) vai MSc (csm tietojenkäsittely) tutkielmaa.
\item Valitse kieli jolla kirjoitat tutkielman: \texttt{finnish}, \texttt{english} tai \texttt{swedish}.
\item Jos olet kirjoittamassa maisterintutkielmaa, valitse linja/opintosuunta.
\end{enumerate}


{\textbf{STEP 2 -- Aseta henkilökohtaiset tietosi}}

\begin{enumerate}
\item Kirjoita alustava otsikko tutkielmallesi: \texttt{\textbackslash title\{\}}.
\item Kirjoita oma nimesi kohtaan \texttt{\textbackslash author\{\}}.
\item Lisää ohjaajien nimet \texttt{\textbackslash supervisors\{\}}.
\item Määrittele avainsanat \texttt{\textbackslash keywords\{\}}.
\item Määritä tutkielmasi ACM luokittelutermit \texttt{\textbackslash classification\{\}}. Ks. lisätietoa: \url{https://dl.acm.org/ccs}.
\end{enumerate}

{\textbf{STEP 3 -- Kirjoita tiivistelmä}}

%\begin{itemize}
%\item 
Voit kirjoittaa tiivistelmän (koko tiivistelmäsivu) eri kielillä \texttt{otherlanguages}-ym\-pä\-ris\-tön avulla. Alla esimerkki jolla kirjoitat englanninkielisen tiivistelmän muulla kuin englannin kielellä kirjoitettuun tutkielmaan:

\begin{verbatim}
\begin{otherlanguage}{english} 
\begin{abstract}
Your abstract text goes here. 
\end{abstract} 
\end{otherlanguage}
\end{verbatim}

%\end{itemize}

{\textbf{STEP 4 -- Kirjoita tutkielma}}

\begin{enumerate}
\item Kirjoittamisesta Latexilla löydät hieman ohjeita alempaa.
\item Poista tämä liite ja muu ohjeistus tutkielmastasi, esim. kommentoimalla.
\end{enumerate}

{\textbf{STEP 5 -- Aseta kirjallisuuslähdeluettelon tyyli}}

\begin{itemize}
\item Oletustyylin tekijä-vuosi, eli (Einstein, 1905), voit vaihtaa viittaustyylin (tiedostossa \texttt{HY-CS-main.tex}) helposti (eri mallit kommentoituna) esim. numeroituun [1], tai aakkostyyliin [Ein05].
Lisää ohjeita liittyen viittaustyylin säätämiseen {Bib}\TeX issä löytyy verkosta: \url{https://ctan.org/pkg/biblatex}
\item Sovi käytettävä tyyli ohjaajasi kanssa. 
\end{itemize}

\section{Kirjallisuusviitteet Latexissa}

Kirjallisuuslähteet ylläpidetään erillisessä .bib-tiedostossa. Tässä tutkielmapohjassa käy\-te\-tyt kirjallisuuslähteet, joista esimerkkejä kuvassa~\ref{bibexamples-fi}, löytyvät tiedostosta\newline \texttt{bibliography.bib}.

\begin{figure}[ht]
    \centering
    \begin{scriptsize}
\begin{verbatim}

@article{einstein,
    author =       "Albert Einstein",
    title =        "{Zur Elektrodynamik bewegter K{\"o}rper}. ({German})
        [{On} the electrodynamics of moving bodies]",
    journal =      "Annalen der Physik",
    volume =       "322",
    number =       "10",
    pages =        "891--921",
    year =         "1905",
    DOI =          "http://dx.doi.org/10.1002/andp.19053221004"
}

@book{latexcompanion,
    author    = "Michel Goossens and Frank Mittelbach and Alexander Samarin",
    title     = "The \LaTeX\ Companion",
    year      = "1993",
    publisher = "Addison-Wesley",
    address   = "Reading, Massachusetts"
}

@book{knuth99,
    author    = "Donald E. Knuth",
    title     = "Digital Typography",
    year      = "1999",
    publisher = "The Center for the Study of Language and Information",
    series    = "CLSI Lecture Notes (78)"
}
\end{verbatim}
\end{scriptsize}
    \caption{Esimerkkejä kirjallisuuslähteiden kuvaamisesta .bib-tiedostossa.}
    \label{bibexamples-fi}
\end{figure}

Viitteet kirjallisuuslähteisiin muodostetaan komennolla \texttt{\textbackslash citep\{einstein\}}, josta generoituu tekstiin valitun viittaustyylin mukaisesti muotoiltu viite \citep{einstein}, tai \texttt{\textbackslash citep\{latexcompanion,knuth99\}}, josta tekstiin puolestaan generoituu \citep{latexcompanion,knuth99}. 
Voit esimerkiksi kirjoittaa \citep{einstein} viitataksesi julkaisuun, jonka on kirjoittanut \citeauthor{einstein} vuonna \citeyear{einstein}, kun vain lähde \citet{einstein} on oikein lisättynä kirjallisuuslähdetiedostossa (katso miltä nämä näyttävät Latex lähdekoodissa).

Tekstissä viitatut kirjallisuuslähteet tulevat automaattisesti viiteluetteloon. Kirjallisuuslähteiden tietojen oikeellisuus ja yhdenmukaisuus .bib-tiedostossa vaikuttavat luonnollisesti siihen, miten tiedot tutkielmassa näyttäytyvät. Tämä on syytä huomioida, sillä esim.\ verkosta valmiiksi {Bib\TeX} muodossa löytyvien tietojen täydellisyyten tai samanmuotoisuuteen ei pidä sokeasti luottaa.  


Keskustele viittaustyylin valinnasta ohjaajan kanssa. 
%Joitain vaihtoehtoja on osoitteessa:\\ 
%\url{https://www.overleaf.com/learn/latex/Biblatex_bibliography_styles}.
%\url{https://www.sharelatex.com/learn/Bibtex_bibliography_styles}.

\section{Joitain ohjeita Latexilla kirjoittamiseen}

Seuraavassa on joitain ohjeita tämän tutkielmapohjan käyttöön maisterintutkielmassa. Kirjoittamisohjeita löytyy useasta eri lähteestä. Voit esimerkiksi tutustua kandidaatintutkielman ohjeisiin. 
Ohjaajan kanssa on hyvä keskustella aikaisessa vaiheessa työn rakenteesta.

Yksityiskohtaisia ohjeita Latexin käyttämäsestä saa parhaiten hakemalla verkosta, esim. haku englanniksi "Overleaf table positioning" tuottaa oletettavasti aika toimivan vastauksen.

\section{Kuvat}
Kuva~\ref{fig:logo-fi} toimii esimerkkinä kuvan lisäämisestä työhön (katso tarkemmin mallia Latex lähdekoodista). Muista myös viitata jokaiseen kuvaan tekstissä. 

\begin{figure}[ht] % remove [h!] for automatic placement, which is probably better for a thesis with more text on page
\centering 
\includegraphics[width=0.3\textwidth]{template/figures/HY-logo-ml.png}
\caption{Helsingin yliopiston logo matemaattis-luonnontieteellisen tiedekunnan värein.\label{fig:logo-fi}}
\end{figure}

\newpage % just to keep the table on the same page with the short piece of text
\section{Taulukot}

Taulukossa~\ref{table:results-fi} on esimerkki kokeellisten tulosten raportoinnista taulukkona. Muista myös viitata jokaiseen taulukkoon tekstissä.
\begin{table}[ht]
\centering
\caption{Kokeelliset tulokset.\label{table:results-fi}}
\begin{tabular}{l||l c r} 
Koe & 1 & 2 & 3 \\ 
\hline \hline 
$A$ & 2.5 & 4.7 & -11 \\
$B$ & 8.0 & -3.7 & 12.6 \\
$A+B$ & 10.5 & 1.0 & 1.6 \\
\hline
%
\end{tabular}
\end{table}



% BSc instructions
%\chapter{Johdanto}


Kaikessa julkaistavaksi tarkoitetussa tekstissä kirjoittajan luomisen ja
esitystavan vapautta rajoittavat monet ohjeet ja tarkatkin määräykset.

Parhaimmillaan lukijalle ja kirjoittajalle yhteinen, tuttu säännöstö luo
eräänlaisen tukiverkoston, joka tukee sanoman siirtymistä vääristymättä.
Kirjoituksen lukija löytää kirjoituksesta helpommin olennaisen sisällön,
jos kirjoituksen ulkoasu ja sisällön rakenne vastaavat hänen
tottumuksiaan. Sama koskee myös kirjoittajaa. Noudattaessaan valmista
esitystapamallia kirjoittajan ei tarvitse käyttää aikaansa itse työn
kannalta toissijaisten seikkojen miettimiseen, vaan hän voi keskittyä
hiomaan tekstin sisältöä. Siksi kannattaa harjoitella myös työn ulkoasua
koskevien ohjeiden noudattamista, vaikka omasta mielestään osaisikin
valita esitykselleen ohjetta paremman muodon.

Tämä kirjoitus on tarkoitettu Helsingin yliopiston
Tietojenkäsittelytieteen osastoon alempien opinnäytteiden ja
harjoitusten ulkoasun ja rakenteen ohjeeksi. Ohje soveltuu siten
kandidaatintutkielman kirjoittamisen kurssille, ohjelmistotuotantoprojekteihin, seminaareihin ja
pro gradu -tutkielmiin. (Kirjoitus on päivitetty uusintapainos aiemmista
ohjeista, jotka kurssin luennoijat ovat laatineet \citep{erkio01,erkiomakela96,erkio94,verkamo92}.)

Tyylimäärittely on saatavissa pdflatex- ja word-versiona.
Tyylimäärittelyitä valitessa on huomattava ohjeet tekstien syöttöön
liittyvässtä koodauksesta (UTF8,ISO 8859-15).
Tämän kirjoituksen tukena sopivat käytettäväksi tavanomaiset latex- tai
word-oppaat. 

\chapter{Kirjoituksen rakenne}

Tarkastellaan aluksi tieteelliseltä tekstiltä odotettuja
kirjoituksen osia. Samoihin asioihin on luonnollisesti syytä
kiinnittää huomiota myös muussa teknisessä kirjoittamisessa. Huomattakoon, että tämä teksti ei ole tieteellinen teksti, eikä siten itse sisällä kaikkia niitä elementtejä, jotka tieteellisen tekstin sisällölliseen antiin kuuluvat. Tällaisia puutteita ovat esimerkiksi johdannon tutkimuskysymyksen asettelun puuttuminen sekä arvoivan materiaalin puute tekstin lopussa, sekä yhteenvedon latteus.
Teksti rajoittuu siten otsikkonsa mukaisesti vain tekniseen sisällön asetteluun.

\section{Tiivistelmä}



Tiivistelmäsivu sisältää seuraavat osat: työn bibliografiset tiedot,
tiivistelmäteksti, aiheluokat ja avainsanat. Bibliografiset tiedot
koostuvat työn otsikosta, tekijän nimestä, julkaisupaikan tiedoista,
julkaisuajankohdasta ja sivumäärästä.

Tiivistelmäteksti on lyhyt, yleensä yhden kappaleen mittainen
(maksimissaan noin 100 sanaa) selvitys
kirjoituksen tärkeimmästä sisällöstä: mitä on tutkittu, miten on
tutkittu ja mitä tuloksia on saatu.


Aiheluokat kuvataan ACM Computing Classification System -luokituksen (CCS)
luokituksen mukaisesti. Luokittelussa käytetään täysia polkuja juurisolmun CCS osoittamista lähtöposteistä lehtisolmuihin. Polkuja voi antaa 1-3 aihepiirien soveltuvuuden mukaan, mitä alempi opinnäyte, sen vähemmän polkuja se tarvitsee. 
Poluissa tasot erotetaan toisistaan nuolella eteenpäin. Kun  polun nimisanoja arvioidaan suhteessa työn sisältöön, merkitään boldface-fontilla tärkein termi, italics-fontilla toiseksi tärkein. Näin menetellään, mikäli jotkin termeistä ovat olennaisesti paremmin kuvaavia kuin muut polun termit. Nimettyjen polkujen lisäksi lukija voi siten tarkastella lisäulottuvuutena myös tärkeiksi merkittyjen termien joukkoa sinänsä.
Avainsanoiksi valitaan kirjoituksen sisältöä
keskeisesti kuvaavia käsitteitä.

\section{Johdanto}


Johdannon tarkoituksena on kertoa yleiskielisesti
työn tavoite. Kerrotaan (kuten tiivistelmässäkin, mutta laveammin),
mitä on tutkittu, miten on tutkittu ja mitä tuloksia on saatu.
Jotta kysymyksenasettelu ja tulokset on lukijan helppo oikein tulkita on syytä aloittaa johdanto asettelemalla tutkimus asiayhteyteensä, esimerkiksi kertomalla aluksi, minkälaisessa yhteydessä tarkasteluun otettavat haasteet esiintyvät ja keiden on ratkaisuista tarkoitus hyötyä.

Johdannon pituus määräytyy suhteessa koko kirjoitelman pituuteen.
Parisivuinen kirjoitus ei erikseen otsikoitua johdantoa kaipaa, sillä
se itsessään on laajennettu tiivistelmä. Kymmensivuisen
kirjoituksen johdanto voi olla vaikkapa sivun tai puolentoista
mittainen. Pro gradu -tutkielman 50-70-sivuiseen kokonaisuuteen
tuntuu 2-4-sivuinen johdanto kohtuulliselta. 

Johdanto kertoo siis lyhyessä, yleistajuisessa muodossa
koko kirjoitelman kysymyksenasettelun, juonen sekä tulokset ja johtopäätelmät.
Tämän luettuaan lukija voi päätellä, haluaako syventyä asiaan tarkemmin
lukemalla koko kirjoituksen.


\section{Käsittelyluvut}

Käsittelylukujen työnjako määräytyy käsiteltävän asian luonteen
mukaisesti.
Lukijan ohjailemiseksi kukin pääluku kannattaa aloittaa lyhyellä
kappaleella, joka paljastaa mikä kyseisen luvun keskeisin sisältö on ja
kuinka aliluvuissa asiaa kehitellään eteenpäin.
Erityisesti kannattaa kiinnittää huomiota siihen, että lukijalle ilmaistaan selkeästi miksi kutakin asiaa käsitellään ja miten käsiteltävät asiat suhtautuvat toisiinsa. 

Jäsentelyongelmista kielivät tilanteet, joissa
alilukuja on vain yksi, tai joissa käytetään useampaa kuin
kahta tasoa (pääluku ja sen aliluvut). Kolmitasoisia
otsikointeja saatetaan tarvita joissakin teknisissä
dokumenteissa perustellusti, mutta nämä muodostavat poikkeuksen.

Perusohjeena on käyttää tekstin rakenteellisesti painokkaita paikkoja,
kuten lukujen avauksia ja teksikappaleiden aloitusvirkkeitä
juonenkuljetukseen ja informaatioaskeleiden sitomiseen toisiinsa.
Tekstikappaleiden keskiosat, samoin kuin lukujen keskiosat selostavat
asiaa vähemmän tuntevalle yksityiskohtia, kun taas aihepiirissä jo
sisällä olevat lukijat voivat alkuvirkkeitä silmäilemällä edetä
tekstissä tehokkaasti eksymättä tarinan juonesta.

Kullakin kirjoittajalla on oma temponsa, joka välittyy lukijalle
tekstikappaleiden pituudessa ja niihin sisällytettyjen ajatuskulkujen
mutkikkuudessa. Kussakin tekstikappaleessa pitäisi pitäytyä vain yhdessä
informaatioaskelessa tai olennaisessa päättelyaskelessa, muuten juonen
seuraaminen käy raskaaksi olennaisten lauseiden etsiskelyksi. Yksivirkkeisiä
tekstikappaleita on syytä varoa. 



\section{Lähdeviittausten käyttö}


Olennaisia opittavia asioita viittaustekniikoissa ovat viitteen paikka
tekstissä, oikea lähdeluettelojärjestys valitun viitetyylin parina sekä
taito ja tahto noudattaa annettua tyylimääräystä. Väitöskirjoissa ja
lehti- tai konferenssiartikkeleissa tekstin hyväksyminen riippuu myös
näiden yksityiskohtien asianmukaisesta käsittelystä. Tästä syystä
laitoksella nähdään tarpeelliseksi opiskelijoiden tutustua edes
pinnallisesti myös muihin tyylilajeihin ja oppia käyttämään
automatisoituja muotoilutyökaluja tehokkaasti, jolloin tyylimuutokset
ovat tehokkaita.

Lähdeviitteet sijoitetaan aina virkkeen sisäpuolelle. Siten esimerkiksi
tekstikappaleen lopussa irrallaan oleva viite ei ole asiallinen. Tilanne
ei muutu, vaikka viite sujautettaisiin tekstikappaleen viimeisen
virkkeen sisään. 
Lähdeviittauksen yhteyteen merkitään mukaan tarkentavat
sivunumerot, mikäli lukijan olisi työlästä löytää asianomainen kohta
viitatusta lähteestä. 


Tehokkaita viitteensijoittelupaikkoja ovat esimerkiksi uuden käsitteen
nimeämiskohta ja virkkeen loppu kun kyseessä on lähteestä lainattu
väite. On myös muistettava lainausmerkkien käyttö silloin kun tehdään
suoria lainauksia.

Tekstin jäsentelyn on tuotava selkeästi esiin, mihin asiaan viite
liittyy. Samalla tulee ymmärrettäväksi se, kuinka pitkään
tekstikatkelmaan ko. viite liitetään. Ei ole siten asiallista aloittaa
lukua nimeämällä yhtä tai useampaa lähdettä luvun taustaksi, vaan
viitteitä on kiinnitettävä täsmällisemmin väitteisiin ja käsitteisiin.
Luvun avaus viitetiedolla voi olla oire myös suuremmasta ongelmasta:
lähderiippuvuudesta. Aloitteleva kirjoittaja helposti toistaa lähteestä
oppimaansa ilman että tarpeellinen analysointi ja prosessointi suhteessa
muuhun opittuun olisi vielä tapahtut.

Viitteillä ja sanamuodolilla on
myös tuotava selkeästi esiin se, mikä teksissä on lainattua ja mikä oman
pohdinnan ja valikoinnin tulosta.


Lähdeviittauksiin käytetään Tietojenkäsittelytieteen osastolla
numeroitua tyyliä ja APA-tyyliä, valinnan näiden välillä tekevät kunkin
ryhmän valvoja ja ohjaaja yhdessä. 
Numeroitu tyyli on esimerkiksi IEEE- ja
ACM-julkaisuissa yleisesti käytetty ja puolustaa siten paikkaansa.
APA-tyyli on poikkeuksellinen ns. kovissa tieteissä, mutta monet
valvojista pitävät siitä sen luettavuuden vuoksi. Numeroita joutuu nimiä
useammin tarkistamaan lähdeluettelosta, sillä tarkastus- ja
arvointiprosessiin kuuluu arvioida myös lähteiden valitaa ja niiden
käyttötapaa.





\section{Yhteenveto}

Yhteenveto  vaatimattomimmillaan on vain lyhyt kertaus kirjoituksen
keskeisistä asioista. Arvokkaamman yhteenvedon saa aikaan kommentoimalla
 työn tulosten arvoa, työn liittymistä ympäristöön ja
tulevaisuudennäkymiä. Tällaiset arviot  huolellisesti
perusteltava.

\section{Lähdeluettelon laatiminen}

Tieteellisen kirjoittamisen kurssin töiden lähdeluetteloiden
laatimisessa noudatetaan seuraavia ohjeita.

Niiden taustalla on kaksi
keskeistä pyrkimystä: tehdä viitatun lähteen hankkiminen luettavaksi
mahdollisimman helpoksi ja ilmaista, millaisen arviointiprosessin
läpi käyneeseen kirjoitukseen vedotaan.
Näistä syistä
\begin{itemize}
\item lähdeviitteen tulee aina olla niin tarkka, että
lähde on sen perusteella tunnistettavissa ja löydettävissä luetteloista
ja kirjastoista,
\item erityyppisten lähteiden (kirjat, konferenssit, lehdet) on erotuttava
toisistaan
ja
\item luettelon eri osien tulee olla mahdollisimman
yhdenmukaisia, erityisesti lähdetyypin sisällä.
\end{itemize}


Riippumatta käytettävästä viitetyylistä, 
lähteet ovat Tietojenkäsittelytieteen osaston opinnäytteiden lähdeluetteloissa tekijän nimen mukaisessa aakkosjärjestyksessä,
saman tekijän (tekijäryhmän) työt julkaisuajan mukaisessa
järjestyksessä. Jos jollakin lähteellä ei ole henkilötekijää, se
aakkostetaan julkaisun nimen mukaisesti. 

Kustakin lähteestä annetaan seuraavat tiedot, edelleen viitetyylistä riippumatta:
\begin{itemize}
\item (tarvittaessa lähdeviitelyhenne).
\item
tekijän tai tekijöiden nimet (sukunimi, etunimien alkukirjaimet)
alkuperäisessä järjestyksessään; jos tekijöitä on enemmän kuin kolme,
voidaan toimia siten, että
vain ensimmäinen tekijä nimetään ja muiden tilalle kirjoitetaan {\em et
al.}
\item
julkaisun tai artikkelin nimi alkuperäisessä muodossaan
\item
julkaisupaikan tiedot:
\begin{itemize}
\item
kirjasta: kustantaja, julkaisupaikka (voidaan jättää pois, jos kyseessä
on tunnettu kustantaja), vuosi ja
\item
lehtiartikkelista: lehden nimi, volyymi, numero, vuosiluku ja kuukausi (suluissa),

\item
artikkelikokoelmassa (esim. konferenssijulkaisussa) ilmestyneestä
artikkelista:
\begin{itemize}
\item kokoelman nimi, toimittaja, kustantaja, julkaisupaikka ja vuosi
{\em tai}
\item konferenssin nimi, järjestäjä, paikka ja aika,
\end{itemize}
\item
raportista: julkaisusarja, raportin numero, julkaisupaikka, julkaisija ja vuosi
ja
\item
www-lähteestä: verkko-osoite, voimassaoloajankohta, mahdollisesti
viittausajankohta hakasuluissa
\end{itemize}
\item
sivunumerot, mikäli lähteenä käytetty julkaisu on artikkeli tai kokoomateoksen itsenäinen luku.
\end{itemize}

Normaaliin suomalaiseen tapaan artikkelin nimessä ainoastaan
ensimmäinen sana kirjoitetaan isolla alkukirjaimella, sen sijaan
konferenssien ja kokoelmajulkaisujen nimissä käytetään isoa
alkukirjainta jokaisen sanan alussa (artikkelisanoja ja prepositioita
lukuunottamatta). Katso mallia oheisista esimerkeistä.
Kokoelman nimen edessä on syytä selvyyden vuoksi käyttää sanaa {\em
Teoksessa}, paitsi kun on kysymys konferenssijulkaisusta, jonka nimi
alkaa lyhenteellä {\em Proc.} (sanasta Proceedings). Tällöin ei tarvita
mitään täydennystä.
Tämän eron näkee esimerkiksi vertaamalla
lähdeviitteiden~''\citep{dantowsley90}''
ja~''\citep{gannonetal89}'' ulkoasuja.

WWW-lähteiden käytössä on syytä muistaa, että verkossa julkaisukynnys on
olematon. Kannattaa siten keskittyä tunnettujen tieteellisten
kustantajien julkaisuihin ja niihin teknisiin standardeihin, joille WWW
on ainoa julkaisukanava. Mikäli sama julkaisu on saatavissa myös
perinteisessä muodossa, viitataan ensisijaisesti siihen ja käytetään
verkko-osoitetta lisätietona. Lähdeluettelossa on annettu esimerkit
useita kanavia julkaistusta kirjoituksesta~\citep{abiteboul,dietinger} sekä pelkästään
WWW-julkaisuna
leviävästä standardista~\citep{bray}.

Erityisesti varoitetaan Wikipedian käytöstä tieteellisessä tekstissä.
Vaikka sen avulla on helppo alustavasti tutustua joihin aihepiireihin ja
asiantuteva lukija voisi teksin kelvolliseksi tiettynä hetkenä
hyväksyäkin, ei se foorumina millään lailla täytä tieteellisesti
vertaisarvoidun tutkimusfoorumin kriteerejä.
Jos Wikipedia-artikkelia  ei mitenkään malta ajankuvana olla mainitsematta, käytettäköön jotain muuta kuin lähdeviitetekniikkaa tähän taiteelliseen otteeseen, vaikkapa alaviitteitä. Olennaista silloinkin on, että tieteellinen sisältö ei tule tällä korvatuksi vaan sen puute korostetuksi.

WWW-lähteeseen viittaamisessa pätevät samat periaatteet kuin
perinteisiin lähteisiin viitattaessa: lähdeviitteessä ilmaistaan
otsakkeet, kirjoittajat, toimittajat js muut seikat. Eroa on ainoastaan
verkko-osoitteen ja sen voimassaoloajankohdan ilmaisemisessa. Mikäli
lähde on julkaistu ainoastaan verkossa, voidaan web-osoitetta (URL)
käyttää vastaavasti kuin perinteisen julkaisun paikannusinformaatiota
(lehden ja se numeron julkaisutiedot). Lähdeluettelossa on WWW-viittausten yhteydessä aina syytä ilmaista päivämäärä, jolloin linkin voimassaolo ja lähteen sisältö on tarkastettu.
Esimerkkeinä verkkoviitteistä soveltuvat seuraavat:
\begin{itemize}
\item Gergen, Kenneth (1999) Narrative, Moral Identity and Historical
Consciousness: a Social Constructionist Account.
http://www.swarthmore.edu/SocSci/kgergen1/text3.html. Haettu 11.6.1999.
\item 	
Ritala-Koskinen, Aino and Valokivi, Heli (2006) The Role of Development
Skills in Social Work Practice Education in Finland. Social Work and
Society, The International Online-Only Journal 4(2006)1.
http://www.socwork.net/2006/1/series/transition/ritalakoskinenvalokivi.
Viitattu 30.8.2006.
\item Heinisuo, Rami and Ekholm, Kai (1997) Elektronisen viittaamisen
opas. Jyväskylän yliopiston kirjaston julkaisuja n:o 40. Jyväskylä:
Jyväskylän yliopiston kirjasto. http://www.pori.tut.fi/~multisil/evo/.
Viitattu 29.8.2006.
\end{itemize}

Kirjoituksen lähdeluettelossa luetellaan täsmälleen ne lähteet, joihin
viitataan kirjoituksen tekstiosassa. Tämän kirjoituksen lähdeluettelo on
tarkoitettu lähinnä esitystavan esimerkiksi, mistä syystä siinä on
''ylimääräisiä'' lähteitä.


Pääsääntöisesti julkaisun tai artikkelin nimen perään tulee piste,
samoin kunkin lähteen bibliografisten tietojen perään. Muut tiedot
erotetaan toisistaan pilkulla. Useimmissa tapauksissa  
voidaan noudattaa teknisten välineiden antamaa mallia, sillä edellytyksellä, että ylläolevat vaatimukset muuten täyttyvät.



\chapter{Ulkoasulliset seikat}

  Tässä luvussa käsitellään yleisimpiä
tekstin tekniseen esittämiseen liittyviä seikkoja.  

\section{Työn osien järjestys}

   Kirjoituksen alussa on aina
erillinen, mallin mukainen kansilehti. Toisena sivuna on
tiivistelmäsivu, sen jälkeen sisällysluettelo (yksi tai useampia
sivuja) ja sitten varsinainen teksti. Sivunumerointi aloitetaan vasta
ensimmäiseltä tekstisivulta (arabialaisella ykkösellä). (Tarkat
jättävät ykkössivun numeromerkittä.) Sisällysluetteloon merkitään
kaikki (numeroidut) otsikot ja vastaavat sivunumerot. Monet
tekstin\-käsittelyjärjestelmät muodostavat itse sisällysluettelon,
jolloin kirjoittajan ei tarvitse huolehtia luettelon sivunumeroiden
päivittämisestä tekstin kehittyessä. Sisällysluettelo\-sivu ja sitä
edeltävät sivut voidaan haluttaessa numeroida erikseen (roomalaisin
numeroin) esimerkiksi tämän mallin mukaisesti.  

Varsinaisen tekstin
jäljessä, mutta itse työhön kuuluvana on ensimmäisenä lähdeluettelo,
jonka otsikkoa ei numeroida. Lähdeluettelon jälkeen sijoitetaan
mahdolliset liitteet, jotka otsikoidaan ja varustetaan sisäisillä
sivunumeroilla.  



Mikäli kuvista, algoritmeista
ja taulukoista halutaan tehdä yhtenäinen luettelo, sijoitetaan
luettelot sisällysluettelon jälkeen. Luetteloiden käyttöarvosta on
eriäviä mielipiteitä, joten niiden laatimiseen ei varsinkaan ilman
tekstinkäsittelyjärjestelmän tukea kannata ryhtyä ilman tarkastajan
erityistä toivetta.  

Mikäli kirjoitukseen erityissyistä halutaan
liittää aakkosellinen hakemisto, sijoitetaan se lähdeluettelon jälkeen
ennen liitteitä. Indeksi merkitään sisällysluetteloon samoin kuin
lähdeluettelo (numeroimaton luku). Mikäli indeksin tekemiseen
ryhdytään, on syytä käyttää tekstinkäsittelyjärjestelmän tarjoamaa
automatiikkaa.

Teksin luonnollisen juonenkuljetuksen mukana esiin
tulevien käsitteiden määrittelyjen sijasta ei pidä yrittää sen enempää
pakata kaikkia määritelmiä johdantoon kuin laatia johdantoa ennen
käsitelistaa tai lyhenteiden selityslistaa. Kumpikaan ei sovi
tavanomaiseen argumentoivaan tieteelliseen tekstityyliin, vaikka
teknisessä yhteydessä niillä liitteinä voi olla lisäarvoa.


\section{Tekstin yleinen sijoittelu}

Lopullinen tutkielmaversio voi olla yksi- tai kaksipuoleiseksi aseteltua
ja riviväliltään 1,5 tai 1.  Erityyppisissä
teksteissä haasteet ja asetteluvaatimukset voivat olla erilaiset. Erota
kappaleet toisistaan yhdellä tyhjällä rivillä tai
käytä tekstinkäsittelytyökalujen ominaisuuksia
hyödyksesi ja määrittele tekstikappaleiden väliin jäävä tila hieman
normaalia riviväliä suuremmaksi.

Kirjoituksen lukujen, kuvien ja taulukoiden erottumisen kannalta
tärkein keino on riittävän tilan käyttö niiden ympärillä. Kuvan ja
nimekkeen tulee olla selkeästi yksi kokonaisuus, joka eroaa muusta
tyhjän tilan rajaamana. Kuvan tai taulukon on aina numerointinsa ja
nimekkeensä kanssa mahduttava yhdelle sivulle tai varmasti
kaksipuolisena paperidokumenttina tarkasteltavassa tekstissä aukeamalle.
Kuvissa fonttikoko ei saa alittaa 8 pistettä.



Jos uusi luku tulisi alkamaan aivan sivun alareunasta (vain yksi tai
kaksi riviä varsinaista tekstiä), aloita mieluummin uusi sivu. Jokaista
uutta lukua ei kuitenkaan ole tarpeen --- etenkään lyhyessä
kirjoituksessa --- aloittaa uudelta sivulta: jos kirjoituksessa on
paljon melkein tyhjiä sivuja, lukija voi epäillä, että kirjoittaja on
yrittänyt saada kirjoituksensa näyttämään pitemmältä kuin se onkaan. 
 Tyhjää tilaa kannattaa käyttää hyödyksi myös kuvien ja taulukoiden
yhteydessä. Erityisesti jos kirjoituksessa käytetään kauttaaltaan samaa
tekstityyppiä, tyhjät rivit ovat välttämättömiä erottamaan esimerkiksi
tekstiä ja taulukkoa toisistaan. Tyhjä tila on halpaa, mutta se lisää
selkeyttä ja luettavuutta.  


\section{Kuvat ja taulukot}


Kuva tai taulukko sijoitetaan mahdollisimman lähelle
(ensimmäistä) tekstikohtaa, jossa siihen viitataan, ei kuitenkaan
kyseistä viittausta aikaisemmaksi.
Tekstissä on syytä myös kertoa, mitä kuvalla halutaan havainnollistaa.
Kuvan voi lukea monella eri tavalla, joten lukijaa on ohjattava.

Kuvaa ei pidä sijoittaa välittömästi luvun otsikon alle, vaan on
aloitettava tekstillä. Kuvaa ei pidä sijoittaa keskelle tekstikappaletta
(saati virkettä), paitsi jos kuva tulee sivun alkuun tai loppuun eikä
kappaleen jatkumisesta tule epäselvyyttä.

Kuvan ei aina tarvitse olla välittömästi viittaavan kappaleen
perässä. Esimerkiksi viittauskohdan ja
vasta seuraavalle sivulle mahtuvan kuvan väliin jäävää sivun loppuosaa
ei jätetä tyhjäksi. Kuvaa ei kuitenkaan pidä viedä seuraavaa
sivua kauemmas viittauskohdasta.


Varsinaista kuvan esittämistä havainnollistaa kuva~\ref{kuvaesimerkki}.
Huomiota on kiinnitettävä kuvan osien ja tekstimerkintöjen näkyvyyteen,
kuvan numerointiin ja otsikointiin. 

\begin{figure}[ht]
%\begin{figure}[tbh] t= top, b = bottom, h=here
\ \newline
\begin{center}
\includegraphics[width=0.75\textwidth]{kuvaesimerkki.pdf}
\caption{Kuvan elementit.}
\label{kuvaesimerkki}
\end{center}
\end{figure}


Kuvien kokoon on kiinnitettävä huomiota. Käytettyjen merkintöjen
on oltava helposti luettavissa ja selkeät. Esimerkiksi
suorituskykykäyriä esitettäessä akselit on nimettävä, asteikot
merkittävä ja käytetyt yksiköt tuotava selkeästi esiin.
Samankaltaisia asioita esitettäessä useammalla kuvalla on
syytä käyttää samaa mittakaavaa vertailun helpottamiseksi.

Kuvan otsikko kirjoitetaan kuvan alle ja sen tulee olla mieluummin lyhyt
ja ytimekäs kuin liian selittelevä.
Samoin toimitaan taulukoiden otsikoinnissa.

Kuvat ja taulukot numeroidaan juoksevasti. Pitkissä teksteissä käytetään
kaksitasoista numerointia (esimerkiksi Kuva 3.1) pääluvuittain, lyhyissä
riittää yksitasoinen numerointi.


Kuva- ja taulukko-otsikoiden yhdenmukaiseen esitystyyliin on syytä kiinnittää
huomiota, samoin mm. välimerkkeihin. Luontevaa on käyttää
kuvatekstin lopussa pistettä, ovathan useimmat kuvateksteistä virkkeitä. 

(Kuvien ja taulukoiden otsikointityyli vaihtelee
kustantajittain ja julkaisuittain. Samoin tuntuu suositeltava käytäntö
Tietojenkäsittelytieteen laitoksen sisällä vaihtelevan taulukon
otsikon sijainnin suhteen.)


\section{Otsikot}

Otsikoissa voi käyttää muusta tekstistä poikkeavaa kirjasintyyppiä,
alleviivausta, suurempaa kirjasinkokoa tms.\ erotuskeinoa, yleensä
kuitenkin vain yhtä näistä, koska kovin monta erilaista kirjasintyyppiä
ja -kokoa tekee ulkoasusta helposti sekavan.  Otsikoiden esitystavan on
oltava johdonmukainen läpi koko kirjoituksen. Numeroimattomia
''ylimääräisiä'' otsikoita ei tule yleensä käyttää.


\section{Mallin käyttö}

Voit käyttää tätä kirjoitusta mallina oman opinnäytteesi ulkoasua
varten. Eri tekstinkäsittelyjärjestelmissä käytössä olevat yksityiskohdat kuten
kirjasintyypit ja -koot ja rivivälit  poikkeavat toisistaan, joten
pienet poikkeamat ovat toki hyväksyttäviä.

Tieteellisen kirjoittamisen kurssin luennoilla ja
liitteenä olevassa ohjeessa annetut töiden ohjeelliset sivumäärät
koskevat työtä, joka vastaa ulkoasultaan tätä ohjetta (kirjasinkoko
12~pistettä). Tässä tekstissä keskimääräinen rivin pituus lienee noin
80~merkkiä ja sivun pituus 35-40~riviä.
Sivumääriin lasketaan varsinaisen tekstiosuuden pituus ja lähdeluettelo
(arabialaisin numeroin numeroitu osuus), ei kansilehteä, tiivistelmää
eikä sisällysluetteloa. Sivumääräarviossa otetaan huomioon hyvin vajaat
sivut, joita syntyy paljon lyhyiden lukujen ja taittotyyliin määritellyn
luvun avauksen pakottaminen oikeanpuolimmaiselle sivulle. 

\chapter{Yhteenveto}

Tämän kirjoituksen tarkoituksena on toimia muistilistana eräistä
esitystavallisista säännöistä, joihin harjoitusten ja tutkielmien
kohdalla on syytä kiinnittää huomiota.

Annetut ohjeet on laitoksen henkilökunta muotoillut yhdessä keskustellen
ja noudattaen oman tieteenalansa perinteitä. Eri erikoistumisaloilla ja
erilaisillaa määräävässä asemassa olevissa julkaisufooruilla käytänteet
vaihtelevat ja nuorten tutkijoiden onkin tiedostettava ero yleisten
sisältöohjeiden ja teknisten muotoilusääntöjen välillä. Aina tekstin
valmistuessa on tarkastettava erikseen, täyttääkö se annetut
pituusrajoitteet ja vastaako se annettuja muotoiluohjeita, olivatpa ne
kuinka pikkutarkkoja tahansa. Tarkasta sääntöjen noudattamisesta syntyy
yhteinäisyyttä kokoovan julkaisun tasolla, mikä helpottaa lukijoiden
työskentelyä.

Tämä ohje vastaa vain asettelullisiin kysymyksiin ja sen rinnalla on
syytä tutustua materiaaliin ja luentoihin, joissa keskitytään tekstin
varsinaiseen sisältöön. Olennaisin väline on kuitenkin akateemisesti
pidemmälle ehtineen, jo julkaisuja rakentaneen ohjaajan palaute ja
mentorointi.


%
\chapter{Introduction}


In all writing for publication, the writer's freedom of creation and expression are limited by a number of guidelines and specific regulations.

At best, a familiar set of regulations shared by reader and writer can create a kind of support network that allows the message to be relayed without distortion. It will be easier for readers to find 
the pertinent contents in a piece of writing if its layout and structure are the same as they are used to. This also applies to writers. When writers follow a set presentation model, 
they do not have to waste time on considerations that are secondary to the work itself, but they can concentrate on polishing the contents of the text. This means that it is a good 
idea to practice following the rules for the layout, though you may think you know how to select a better way to present your work.

This is a guide for the layout and structure of theses and essays at the Department of Computer Science at the University of Helsinki. It is thus applicable to the course 
Scientific Writing, the software engineering projects, seminars, and MSc theses. (This is an updated version of 
the previous guide written by the course lecturers \citep{erkio01, erkiomakela96, erkio94, verkamo92}.)

The \LaTeX\ guide and \LaTeX\ style that has been  published on the department's web site can be used as support for this guide.


\chapter{Structure}

Let us start by looking at the sections expected to be in a scientific text. Keep in mind that the same expectations go for all kinds of technical writing.
However, this document in itself, is not a scientific or a research
text, so there will be content lacks in terms of research question
setting in the introduction and evaluative material in the last sections.

\section{Abstract or summary}
%\enlargethispage{5mm}


The summary page contains the following elements: the bibliographical data of the work, an abstract, topic classification, 
and the key words. The bibliographical data consists of title, name of the author, place of publication, date of publication, and number of pages.

The abstract should be short, generally one paragraph (100 words maximum) explaining the main contents of the work: topic, methodology and results.


 Topics are classified according to the ACM Computing Classification System
(CCS). A small set of paths (1-3) should be used, starting from any top nodes
referred to bu the root term CCS leading to the leaf nodes. The elements
in the path are separated by right arrow, and emphasis of each element individually can be indicated
by the use of bold face for high importance or italics for intermediate
level. The combination of individual boldface terms may give the reader
additional insight. 

\section{Introduction}


The purpose of the introduction is to introduce the goals of the work in general terms. Describe topic, 
methodology and results (as in the abstract, but expand it).
In order to provide the reader a good starting point for
interpretations, it is good to start the introduction by
contextualisation of the challenges and solutions to be discussed. For
example, why a certain domain has a particular challenge and who are
intended to benefit from the solutions proposed.

The length of the introduction depends on the length of the whole work. A few pages of text does not need a 
separate introduction, since it is an expanded summary in itself. The introduction to a 10-page text can 
be 1--1.5 pages long. For a 50--70-page MSc thesis, a 2--4-page introduction seems reasonable.

The introduction should shortly describe the problem field of the whole work, the plot, and the conclusions, 
in general terms. After reading it, the reader may decide whether to go deeper into the topic by reading the whole text.


\section{Topic chapters}

The nature of the matter at hand determines how the topic chapters are disposed.
In order to guide the reader, it is a good idea to start each main chapter with a short paragraph 
on what the main topic of the chapter is and how it progresses from one sub-chapter to the
next. Especially relevant is to express how concepts, challenges,
solutions and research steps are bound together. There should be enough
guidance for the reader to allow expectation of the right storyline.


Basic rule for easy to follow text is to use the natural emphasis of the
text structure to support the content matter key concepts and thought
processes. This means using openings of sectiosn and text paragraphs for
key arguments and information moves, while the internal parts of
paragraps are filled with supporting aspects that those less familiar
with the topic area need. Those with better background knoweldge can
quickly skim trhough the text without loosing any essential arguments.

Each author has his or her personal ryhtm in the text, which is visible
for the reader as the length of text paragraphs and complexity of
thought chains within. A good policy is to take only one information
move or transition per paragraph. This way the text stays easy to follow.

Signs of problems with the disposition of the text are easily seen in texts
with only one sub-chapter, or with more than two chapter levels (main and sub-chapters). There may be justifiable reasons to use three-level 
headings in some technical documents. Single sentence text paragraphs
are also to be avoided.



%\pagebreak
\section{Reference usage}


Relevant learning targets include superficial knowledge of several
citation styles and capability (and willingness) to follow a given
style and ordering of entires and bibliographical details in the list of references.
These aspects are essential as the approval of a PhD manuscript or
journal article may depend on them.

Disregard of which style you use, 
references are always placed inside sentences.  This means that e.g. a separate reference at the end of a paragraph would be inappropriate.

The structure of the text must clearly show what the reference relates to.  At the same time, it 
shows how long a piece of the text that the reference relates to.

Efficient positions for citations are right after the introduction
(definition) of a concept, a methods or such, or in the end of a claim
from the reference material. Furthermore, if quoting verbatime one must
use citation marks.

The text structures and wordings, in addition to the location of
ciations must clearly express whether claims or arguments are 
authors' own, or if they come from the contribution of others. Thus is
is of bad style to open a section by listing the citations on which the
section is based on. Such method can further indicate more serious
problems like following reference material as it is instead of analysing
and synthetising material into new though processes. 


\section{Conclusion}

At its simplest, a conclusion is merely a weak revision of the main points of the text.  All more valuable 
conclusion sections contain comments on e.g. the value of results, how the work relates to its environment, or 
future visions. This kind of evaluations should be well-grounded in fact, though, or the conclusion 
might inadvertently seem comical. 

\section{Creating the list of references}

The following guidelines should be followed when creating lists of reference for the 
assignments during the course Scientific Writing.

The guidelines are backed by two main goals: to make it as easy as possible to find the 
referenced source, and to show what kind of evaluation process the referenced work has undergone. For these reasons
\begin{itemize}

\item the reference notes should always be so exact that the source can be recognized and found in catalogues and libraries

\item different types of sources (monographs, conferences, journals) have to be easy to distinguish from each other, and 

\item the different parts of the list must conform to each other, especially for each source type.
\end{itemize}


Independent of the citation style in use, 
the sources are listed alphabetically according to the author's name, and works by the same 
author (group of authors) in the order in which they have been published. If some publication 
does not have an individual writer, it is alphabetized according to its name.


The following information should be given on each source independent of
the citation style:
\begin{itemize}
\item
Name(s) of author(s) (surname, initial letters of first names) in the original order; if there are more 
than three authors, we can write the name of the first author and {\em et al.} instead of the other names.

\item
the name of the publication or article in its original form

\item
place of publication:

\begin{itemize}

\item
of monographs: publisher, place of publication (can be omitted if the publisher is well known), year
\item
of journal articles: name of journal, volume, issue, year and month (in parenthesis)

\item
of articles from article collections (such as conference publications):
\begin{itemize}

\item name of collection, editor, publisher, place of publication and year 

{\em or}

\item conference name, coordinator, place and time,

\end{itemize}

\item
of a report: series, report number, place, publisher and year

\item
of a web source: URL, validity date, possibly the date when referenced in square brackets

\end{itemize}

\item
page numbers, if the source is an article or constitutes a chapter in a compilation.
\end{itemize}


When using web sources, you should keep in mind that the threshold for publication on the web 
is non-existent. It is better to concentrate on the publications of well-known scientific 
publishers and the technical standards for which the web is the only publication channel. If 
the same publication is available on paper, refer to that primarily and add the URL as a complement.


The list of references gives an example of a text that has been published through many 
channels~\citep{abiteboul,dietinger}and another example that shows a standard that has 
been disseminated only through the web~\citep{bray}. For web-based
references it is important to give both the publication date and the
time of reading and interpreting the material. The content is prone to
change, and dating your use, you protect your interpretations for
unnecessary accusations of being faulty, if it is known which version
you used.



The list of references for a text should list exactly the sources that the text refers to. 
The list of references for this text is an example of how to present sources, 
which is why it contains ''extra'' sources.


In any style,
put a full stop after the name of a publication or article, as well as after the bibliographical 
data of each reference. Separate the other pieces of information with a comma. As is normal in 
Finnish, only the first letter of the first word in the heading is capitalized, but in the titles of 
conferences and compilations, each major word is capitalized (not articles or prepositions). See the 
appended example for a model. For the sake of clarity, it is best to write {\em In the work} before 
the name of a compilation, except in the case of conference publications where the name starts with 
the abbreviation {\em Proc.} (for Proceedings). In such cases, no complement is necessary. You can 
see the difference by comparing the layouts of the references ~''\citep{dantowsley90}'' and~''\citep{gannonetal89}'' .
In case of other citation styles, it is likely that you can trust the
results given by the automated reference management tools.

\chapter{Layout}

This chapter discusses the main issues of the technical presentation of a text.


\section{Disposition of text sections}


Each text should include a separate cover page, as in this model. The second page contains an 
abstract, followed by a table of contents (on one or more pages), and then the main text.  
Pagination starts on the first page of the main text (with the Arabic numeral 1). (Rigorous writers 
leave out the number from the first page.) All the (numbered) headings and their page numbers should be 
written into the table of contents. Many word-processing systems create the table of contents automatically 
so that the writer does not have to worry about updating page numbers as the work progresses. If writers 
wish to, they can paginate the table of contents and previous pages separately (with Roman numerals), e.g. like this model.

After the main text, but as part of the text body, comes the list of references; its heading is 
not numbered. Any possible appendices are added after the list of references, with headings and internal page numbers.

%\enlargethispage{5mm}

If you want to make a coherent list of figures, algorithms and tables, this list should be placed i
mmediately after the table of contents. The value of such lists is a matter of opinion, so there 
is no need to create one --- especially if your word-processing software does not support it --- unless your instructor explicitly asks for it.

If for some special reason you want to add an alphabetical index top your work, it should be placed after 
the list of references and before appendices. An index should be noted in the table of contents, as s
hould the list of references (unnumbered chapter). Writers who want to create an index should use the automatic support their word-processing system offes.

\section{General text layout}

Nowadays you may plan your text to be printed on single side or double
side, with line spacing one, for the final copy.

For reviewing and feedback purposes, check with your readers for their
preferences. Most read and comment on paper and need some space for
marking. 
For drafts you can use wide margins too. 
Both margins should also be fairly broad (ca.~3~cm): the wider left margin is needed for binding the text and the
right one for the instructor's comments.  Leave enough space (2--3~cm) at the top and bottom, as well.

The most effective way to distinguish features like chapters, figures, etc,  is to have enough space 
in your text.  Separate paragraphs from each other with one and chapters with a few empty lines.  If a 
new chapter is about to start at the bottom of a page (with only one or two lines of text), it is better 
to start it on the next page.  However, it is not necessary to start every chapter on a new page, especially 
not with a short text; if the text contains many pages that are nearly empty, readers might suspect that the 
writer has tried to make it look longer than it is.


Any empty space may be utilized for displaying figures and tables.  Especially if a text is written with the 
same type of text throughout, empty lines are necessary for separating e.g.~text from tables.  Empty space 
is cheap, but adds to the clarity and readability of a text.

\section{Figures and tables}

%\enlargethispage{5mm}

All figures and tables should be placed as near the (first) place in the text that refers to them, but 
not before this reference. The text should also explain what the writer wants to illustrate with the 
figure. Figures can be interpreted in many different ways, so the reader needs guiding.

Figures should never be placed immediately under the heading of a chapter, but chapters should always 
start with text. Figures should not be placed in the middle of a paragraph (much less a sentence), except 
if the figure is placed at the top or bottom of a page and it is clear that the paragraph continues.

Figures do not always have to be placed immediately after the paragraph that refers to them. If there is 
not enough space for a figure on the same page as the paragraph that refers to it, the rest of the page 
should not remain empty, Though the figure is inserted on the following page. However, the figure should 
never be more than one page away from the reference.


The image~\ref{kuvaesimerkki} shows how to present a figure. You must pay attention to the visibility of 
figure parts and text, to the numbering of figures, and captions. 

\begin{figure}[ht]
%\begin{figure}[tbh] t= top, b = bottom, h=here
\ \newline
\begin{center}
\includegraphics[width=0.9\textwidth]{kuvaesimerkki.pdf}
%\rotatebox{90}{\includegraphics[scale=.75]{kuvaesimerkki.pdf}}
\caption{Figure elements.}
\label{kuvaesimerkki}
\end{center}
\end{figure}


We must also pay attention to the size of figures. Any annotations must be clear and easy to read When 
presenting performance curves, for example, the axels have to be named, the scale notated, and the units 
clearly presented. If you present many things with similar figures, you should use the same scale for easy comparison.

The caption of a figure should be written underneath it, and it is preferable that it be short and to 
the point than too explicit. The same goes for table captions.

Figures and tables should be numbered progressively. In long texts, two-level numbering should be 
used (e.g. Figure 3.1) according to the chapter number, but in shorter texts, one-level numbering is good enough.


You should pay attention to presenting figure and table captions consistently, as well as to punctuation 
marks. It is natural to put a full stop after a caption, since they are most often full sentences. 

(The style for figure and table captions vary according to publisher and publication. The recommendations 
at the Department of Computer Science also seems to vary with respect to where the caption should be placed.)


\section{Headings}

You can use a different font in headings than in the rest of the text, or underlining, larger fonts, or other 
methods of emphasis, but generally it is preferred that you only use one of these methods, because if there are 
very many font types and sizes, the layout looks messy.  The format of headings must be consistent throughout 
the text. You should not use any unnumbered ''extra'' headings.


\section{Using this model}

You can use this text as a model for the layout of your own assignment. The font types and sizes, 
line spacing, etc vary according to word-processing system, so small deviations from the rule are acceptable.

The directive number of pages for written assignments given during the lectures for the Scientific-writing 
course and in this guide are applicable to assignments with a layout like this guide's (font size 12 points). The 
average line in this text should consist of about 80~characters and one page about 30~lines. The number of pages 
includes the text itself and the list of references (the part paginated with Arabic numbers), not the 
cover page, summary, or table of contents.

\chapter{Conclusion}

This text is a checklist for some of the rules governing written presentations, which you should 
keep in mind when writing exercises and theses.

This collection of advise has been compiled by staff members as the
result of discussions on general good writing habbits in computer
science. The consensus is that all young researchers and academic degree
holders must be able to seek and follow instructions of layout and
structuring of texts depending on the contribution they are working on. 
This set of instructions aims at unifying the looks of theses and other
study reports from the department, but it is also a representative of
the instruction sets students will meet later.

This set of instructions only address the looks of the document. Other
instructions must be sought, read and gained experience with in order to
learn to select the suitable semantical contents for the scientific texts.



\end{appendices}
%%%%%%%%%%%%%%%%%%%%%%%%%%%%%%%%%%%%%%%%%%%%%%%%%%%%%%%%%

\end{document}
